% =====================================================================
%  A Gentle Introduction to Inverse Reinforcement Learning
%  What Must Be Assumed, and What Must Be Approximated
%
%  Yuhan Chi -- Fudan University
%  August 19, 2026
% =====================================================================
\documentclass[11pt]{article}

% ---------------------------------------------------------------
% Page geometry
% ---------------------------------------------------------------
\usepackage[a4paper,margin=1in]{geometry}

% ---------------------------------------------------------------
% Math
% ---------------------------------------------------------------
\usepackage{amsmath,amssymb,amsfonts,mathtools}
\usepackage{amsthm}

% ---------------------------------------------------------------
% Fonts / encoding
% ---------------------------------------------------------------
\usepackage[T1]{fontenc}
\usepackage[utf8]{inputenc}

% ---------------------------------------------------------------
% Tables, lists, misc layout
% ---------------------------------------------------------------
\usepackage{booktabs}
\usepackage{array}
\usepackage{enumitem}
\usepackage{longtable}
\usepackage{multirow}

% ---------------------------------------------------------------
% Colour and boxes
% ---------------------------------------------------------------
\usepackage{xcolor}
\usepackage[most]{tcolorbox}
\tcbuselibrary{breakable,skins}

% ---------------------------------------------------------------
% Algorithms
% ---------------------------------------------------------------
\usepackage{algorithm}
\usepackage{algpseudocode}

% ---------------------------------------------------------------
% Links (loaded late, as usual)
% ---------------------------------------------------------------
\usepackage{hyperref}
\hypersetup{
  colorlinks   = true,
  linkcolor    = noteblue,
  citecolor    = noteblue,
  urlcolor     = noteblue,
  linktoc      = all,
  pdftitle     = {A Gentle Introduction to Inverse Reinforcement Learning},
  pdfauthor    = {Yuhan Chi},
}

% =====================================================================
% Colours
% =====================================================================
\definecolor{noteblue}{RGB}{20,45,140}
\definecolor{grayhead}{RGB}{80,80,80}
\definecolor{graybody}{RGB}{246,246,246}
\definecolor{redhead}{RGB}{120,15,15}
\definecolor{redbody}{RGB}{253,242,242}
\definecolor{navyhead}{RGB}{15,20,90}
\definecolor{navybody}{RGB}{240,241,250}

% =====================================================================
% tcolorbox styles
%   graybox   -- neutral / expository call-out boxes
%   redbox    -- warnings, common errors, "what this means"
%   navybox   -- Ledger boxes and definition-style displayed boxes
%   navyplain -- untitled continuation box (same look as navybox)
% =====================================================================
\newtcolorbox{graybox}[1]{%
  enhanced, breakable, sharp corners=downhill,
  colback=graybody, colframe=grayhead!70!black,
  colbacktitle=grayhead, coltitle=white,
  fonttitle=\bfseries, boxrule=0.5pt, top=2mm, bottom=2mm,
  title={#1}
}

\newtcolorbox{redbox}[1]{%
  enhanced, breakable, sharp corners=downhill,
  colback=redbody, colframe=redhead,
  colbacktitle=redhead, coltitle=white,
  fonttitle=\bfseries, boxrule=0.5pt, top=2mm, bottom=2mm,
  title={#1}
}

\newtcolorbox{navybox}[1]{%
  enhanced, breakable, sharp corners=downhill,
  colback=navybody, colframe=navyhead,
  colbacktitle=navyhead, coltitle=white,
  fonttitle=\bfseries, boxrule=0.5pt, top=2mm, bottom=2mm,
  title={#1}
}

% untitled box, navy border -- used for "Delivers." style paragraphs
\newtcolorbox{navyplain}{%
  enhanced, breakable, sharp corners=downhill,
  colback=navybody, colframe=navyhead,
  boxrule=0.5pt, top=2mm, bottom=2mm,
}

% =====================================================================
% Theorem-like environments
%   plain style (italic body): theorem, lemma, proposition, corollary
%   definition style (upright body): definition, remark, example, assumption
% =====================================================================
\theoremstyle{plain}
\newtheorem{theorem}{Theorem}[section]
\newtheorem{lemma}[theorem]{Lemma}
\newtheorem{proposition}[theorem]{Proposition}
\newtheorem{corollary}[theorem]{Corollary}

\theoremstyle{definition}
\newtheorem{definition}[theorem]{Definition}
\newtheorem{assumption}[theorem]{Assumption}
\newtheorem{example}[theorem]{Example}
\theoremstyle{remark}
\newtheorem{remark}[theorem]{Remark}

% =====================================================================
% Convenience macros (matching the notation fixed in the handout)
% =====================================================================
\renewcommand{\S}{\mathcal{S}}
\newcommand{\A}{\mathcal{A}}
\newcommand{\T}{\mathcal{T}}
\renewcommand{\P}{P}
\newcommand{\E}{\mathbb{E}}
\newcommand{\R}{\mathbb{R}}
\newcommand{\Prb}{\mathrm{Pr}}
\newcommand{\KL}{\mathrm{KL}}
\newcommand{\JS}{\mathrm{JS}}
\newcommand{\dif}{\,d}
\DeclareMathOperator*{\argmax}{arg\,max}
\DeclareMathOperator*{\argmin}{arg\,min}
\newcommand{\soft}{\mathrm{soft}}
\newcommand{\samp}{\mathrm{samp}}
\newcommand{\GA}{\mathrm{GA}}

% =====================================================================
% Title-page data
% =====================================================================
\title{\vspace{-1.5em}\textbf{A Gentle Introduction to Inverse Reinforcement
Learning}\\[0.3em]
\large\textit{What Must Be Assumed, and What Must Be Approximated}}
\author{\textbf{Yuhan Chi}\\[0.2em]\normalsize Fudan University\\[0.2em]
\href{https://github.com/chi-shan0707}{\texttt{https://github.com/chi-shan0707}}}
\date{August 19, 2026}

\begin{document}

\maketitle
\thispagestyle{empty}

\begin{abstract}
\noindent
Forward reinforcement learning is a well-posed problem: a reward determines
an optimal value function uniquely, and a contraction mapping computes it.
Inverse reinforcement learning is not well-posed, and it fails in two
independent ways at once. It fails \emph{statistically}, because infinitely
many rewards explain the same behaviour, so no amount of data can single one
out without extra assumptions. And it fails \emph{computationally}, because
every principled way of scoring a candidate reward requires an expectation
under the policy that reward induces---an entire forward RL problem nested
inside every gradient step.

These notes use that pair as their spine. Each method is presented as a
specific answer to two separate questions: \emph{what did you assume in
order to make the answer unique?} and \emph{how did you compute the inner
expectation?} Read this way, the sequence Max-Margin $\to$ MaxEnt $\to$ GAIL
$\to$ AIRL stops looking like a list of tricks and becomes a single
argument.

We develop Maximum-Entropy IRL first in the finite, known-model setting,
where the partition function is an honest finite sum and can be computed
exactly by a soft Bellman backward pass. Only after the reader can point at
the two places where $P(s' \mid s,a)$ is consumed do we ask what to do
without it. Part~\ref{part:exits} then treats, in order, the three
questions that trouble almost every beginner: what changes when the problem
stops being finite, whether the reward features must be hand-crafted, and
what remains possible when the transition kernel is unknown. We are
deliberately careful about four points that are frequently stated
incorrectly in secondary sources: the maximum-entropy objective requires a
dynamics-induced reference measure from its very first line; naive maximum
entropy is risk-seeking under stochastic dynamics and must be replaced by
maximum \emph{causal} entropy; GAIL's discriminator carries \emph{no} reward
information at convergence; and AIRL's transfer guarantee requires a
\emph{state-only} reward term.
\end{abstract}

\tableofcontents

\clearpage

% =====================================================================
\section*{How to Read These Notes}
\addcontentsline{toc}{section}{How to Read These Notes}

A first course in IRL usually presents four algorithms in historical order
and leaves the reader to infer why each one replaced the last. These notes
instead fix two coordinates and locate every method on them.

\begin{graybox}{The Two Coordinates}
\textbf{Coordinate 1 --- The identification problem: \emph{what must be
assumed?}}
Behaviour does not determine reward. The forward map $R \mapsto \pi^\star$
destroys information, and no estimator can recreate it. Every IRL method
therefore smuggles in (i) a \emph{behaviour model} saying how reward
produces action, and (ii) a \emph{selection principle} choosing one
representative from the surviving equivalence class, and (iii) a
\emph{reward class} restricting what $R$ may look like. These are modelling
commitments, not discoveries. A method is only as trustworthy as they are.

\medskip
\textbf{Coordinate 2 --- The computation problem: \emph{what must be
approximated?}}
Every selection principle scores a candidate reward $R$ by comparing the
expert's behaviour statistics against those of the policy that $R$ induces.
Producing that second quantity means solving a forward RL problem
\emph{inside} the loop, and then taking an expectation under its solution.
In a finite known MDP both steps are exact. Outside that setting they are
not, and the history of IRL is largely the history of progressively
cheaper surrogates for that one expectation.
\end{graybox}

The two coordinates are genuinely independent, and keeping them apart
resolves most early confusion. A method can be statistically honest and
computationally hopeless (exact tabular MaxEnt in a large MDP), or
computationally excellent and statistically silent about reward (GAIL).
When someone says ``IRL is ill-posed,'' ask which coordinate they mean.

\paragraph{The route}
\begin{itemize}
  \item Part~\ref{part:wellposed} establishes the forward direction and,
  more importantly, isolates the four properties that make it well-posed,
  so that we can watch each one fail in turn.
  \item Part~\ref{part:illposed} attacks Coordinate 1 head-on: what
  precisely can and cannot be inverted, and what the honest target of IRL
  therefore is.
  \item Part~\ref{part:selection} develops Max-Margin and Maximum-Entropy
  IRL in the one regime where Coordinate 2 costs nothing---finite spaces,
  known $P$. This is the reference point for everything after it. The
  partition function is a finite sum; we compute it two ways and check the
  numbers agree.
  \item Part~\ref{part:exits} leaves that regime along its three natural
  exits, which are exactly the three questions beginners ask: finite
  $\to$ infinite, given $\to$ learned features, model-based $\to$
  model-free.
  \item Part~\ref{part:adversarial} shows GAIL and AIRL are the answers
  those three exits force on us, not an unrelated import from generative
  modelling.
  \item Part~\ref{part:synthesis} puts every method back on the two
  coordinates in a single table.
\end{itemize}

Each method section closes with a \textbf{Ledger} box recording, in the
method's own terms, what it assumed and what it approximated. If you
remember nothing else from these notes, remember to fill in that ledger
for any IRL paper you read.

\clearpage
% =====================================================================
\section*{Notation}
\addcontentsline{toc}{section}{Notation}

Symbols are reused across the IRL literature with maddening inconsistency.
We fix a single convention here and keep it for the whole document. In
particular: $P$ always denotes the transition kernel and never a
distribution over trajectories (those are lower-case $p$); the discount is
always $\gamma$ and never $\lambda$; $\phi$ is always a per-step feature and
$\Phi$ always its discounted trajectory sum; the entropy temperature is
$\beta$; and the margin variable is $\xi$, so that $t$ is free to mean time
and $k$ to mean an iteration index.

\begingroup
\small
\setlength{\tabcolsep}{4pt}
\renewcommand{\arraystretch}{0.95}
\begin{center}
\begin{longtable}{@{}p{0.22\linewidth}p{0.38\linewidth}p{0.32\linewidth}@{}}
\multicolumn{3}{@{}l}{\textit{Table 1: Notation used throughout.}}\\[2pt]
\toprule
\textbf{Symbol} & \textbf{Meaning} & \textbf{Type / definition}\\
\midrule
\endhead
\multicolumn{3}{@{}l}{\textit{The environment}}\\
$\S,\A$ & State and action spaces & finite unless stated otherwise\\
$P(s'\mid s,a)$ & Transition kernel & $\S\times\A \to \Delta(\S)$\\
$\mathbf{P}^\pi$ & State-to-state matrix induced by $\pi$ & $\in \R^{|\S|\times|\S|}$, row-stochastic\\
$d_0(s)$ & Initial state distribution & $\in \Delta(\S)$\\
$\gamma$ & Discount factor & $\gamma \in (0,1)$\\
$T$ & Horizon (finite-horizon sections only) & $T \in \mathbb{N}$\\
$\mathcal{M}$ & MDP without reward, $(\S,\A,P,\gamma,d_0)$ & the ``dynamics''\\
\addlinespace
\multicolumn{3}{@{}l}{\textit{Reward, policy, value}}\\
$R(s,a)$ & Reward function & $\S\times\A \to \R$, $|R|\le R_{\max}$\\
$R_w(s,a)$ & Reward in a parametric class & e.g. $R_w(s,a)=w^\top\phi(s,a)$\\
$\pi(a\mid s)$ & Stationary policy & $\S \to \Delta(\A)$\\
$\pi_E$ & Expert policy & assumed (near-)optimal for $R^\star$\\
$V^\pi, Q^\pi$ & Value functions under $\pi$ & $\S\to\R$, $\S\times\A\to\R$\\
$V^\star, Q^\star$ & Optimal value functions & fixed points of $\mathcal{B}^\star$\\
$V_\soft, Q_\soft$ & Soft (entropy-regularised) value functions & fixed points of $\mathcal{B}_\soft$\\
$\mathcal{B}^\star, \mathcal{B}_\soft$ & Bellman optimality / soft operators & $\R^{|\S|}\to\R^{|\S|}$\\
$J_R(\pi)$ & Expected discounted return of $\pi$ under $R$ & $\E_\pi[\sum_t \gamma^t R(s_t,a_t)]$\\
$\beta$ & Entropy temperature / regularisation weight & $\beta > 0$\\
$\Psi(s)$ & Potential function for reward shaping & $\S\to\R$\\
\addlinespace
\multicolumn{3}{@{}l}{\textit{Trajectories and data}}\\
$\tau$ & Trajectory & element of $\T$\\
$\T$ & Set of feasible trajectories & finite when $\S,\A,T$ finite\\
$\mathcal{D}_E$ & Expert dataset $\{\tau^{(1)},\ldots,\tau^{(m)}\}$ & finite demonstrations\\
$p_\pi(\tau)$ & Trajectory distribution induced by $\pi$ & $d_0 \prod_t \pi P$\\
$q(\tau)$ & Dynamics-induced reference measure & $d_0(s_0)\prod_t P(s_{t+1}\mid s_t,a_t)$\\
\addlinespace
\multicolumn{3}{@{}l}{\textit{Features and occupancy}}\\
$\phi(s,a)$ & Per-step feature map & $\S\times\A\to\R^d$\\
$\Phi(\tau)$ & Discounted trajectory feature & $\sum_{t\ge0}\gamma^t\phi(s_t,a_t)\in\R^d$\\
$\mu(\pi)$ & Feature expectation of $\pi$ & $\E_{\tau\sim p_\pi}[\Phi(\tau)]$\\
$\mu_E$ & Empirical expert feature expectation & $\frac1m\sum_i \Phi(\tau^{(i)})$\\
$\rho_\pi(s,a)$ & Discounted occupancy measure & $(1-\gamma)\sum_t \gamma^t \Prb_\pi(s_t=s,a_t=a)$\\
$\rho_E$ & Expert occupancy measure & estimated from $\mathcal{D}_E$\\
$w$ & Reward weight vector & $w\in\R^d$\\
\addlinespace
\multicolumn{3}{@{}l}{\textit{Maximum entropy}}\\
$\mathcal{H}_q(p)$ & Entropy of $p$ relative to $q$ & $-\sum_\tau p(\tau)\log\frac{p(\tau)}{q(\tau)}$\\
$Z(w)$ & Partition function & $\sum_{\tau\in\T} q(\tau)e^{w^\top\Phi(\tau)}$\\
$c$ & Lagrange multiplier for normalisation & $c\in\R$\\
$\xi$ & Max-margin slack/margin variable & $\xi \in \R$\\
\addlinespace
\multicolumn{3}{@{}l}{\textit{Adversarial methods}}\\
$D(s,a)$ & Discriminator; $\Prb[$pair came from expert$]$ & $\S\times\A\to(0,1)$\\
$g_\theta(s)$ & AIRL task-reward network (state-only) & $\S\to\R$\\
$h_\omega(s)$ & AIRL shaping network & $\S\to\R$\\
$f_{\theta,\omega}$ & AIRL discriminator logit & $g_\theta(s)+\gamma h_\omega(s')-h_\omega(s)$\\
$\psi,\psi^*$ & Reward regulariser and its convex conjugate & $\psi:\R^{\S\times\A}\to\overline{\R}$\\
\bottomrule
\end{longtable}
\end{center}
\endgroup

\clearpage

% =====================================================================
\part{The Well-Posed Direction}\label{part:wellposed}
% =====================================================================

\section{Forward Reinforcement Learning}

Our purpose in this part is not to teach RL from scratch but to isolate
\emph{why} the forward problem is easy, in a form precise enough that we
can later watch each reason fail.

\subsection{Markov decision processes}

\begin{definition}[Markov Decision Process]
A finite MDP is a tuple $(\S,\A,P,R,\gamma,d_0)$ where $\S$ and $\A$ are
finite sets, $P(s'\mid s,a)$ is a transition kernel, $R:\S\times\A\to\R$ is
bounded by $R_{\max}$, $\gamma\in(0,1)$, and $d_0\in\Delta(\S)$. We write
$\mathcal{M}$ for the reward-free part $(\S,\A,P,\gamma,d_0)$ and call it
the \emph{dynamics}.
\end{definition}

The separation of $\mathcal{M}$ from $R$ is not cosmetic. IRL is the
problem of recovering $R$ \emph{given} $\mathcal{M}$ (or given the ability
to interact with it) and given behaviour. Everything hinges on how much of
$\mathcal{M}$ we are allowed to know.

For a stationary policy $\pi$,
\begin{equation}
V^\pi(s) \triangleq \E_\pi\Big[\sum_{t=0}^\infty \gamma^t R(s_t,a_t) \,\Big|\, s_0=s\Big],
\qquad
Q^\pi(s,a) \triangleq R(s,a) + \gamma\sum_{s'\in\S} P(s'\mid s,a)V^\pi(s'),
\label{eq:vqdef}
\end{equation}
and the scalar objective is $J_R(\pi) \triangleq \E_{s_0\sim d_0}[V^\pi(s_0)]$.

\subsection{Policy evaluation is a linear solve}

Fix $\pi$. Define $\mathbf{V}^\pi \in \R^n$ with $n=|\S|$, the induced
reward vector $[\mathbf{R}^\pi]_i = \sum_a \pi(a\mid s_i)R(s_i,a)$, and the
induced transition matrix $[\mathbf{P}^\pi]_{ij} = \sum_a \pi(a\mid
s_i)P(s_j\mid s_i,a)$. The Bellman expectation equation is then linear:
\begin{equation}
\mathbf{V}^\pi = \mathbf{R}^\pi + \gamma \mathbf{P}^\pi \mathbf{V}^\pi.
\label{eq:bellmanexp}
\end{equation}

\begin{theorem}[Exact policy evaluation]\label{thm:policyeval}
For $\gamma\in(0,1)$ and row-stochastic $\mathbf{P}^\pi$, the matrix
$\mathbf{I}-\gamma\mathbf{P}^\pi$ is invertible and
\begin{equation}
\mathbf{V}^\pi = (\mathbf{I}-\gamma\mathbf{P}^\pi)^{-1}\mathbf{R}^\pi
= \sum_{k=0}^\infty \gamma^k (\mathbf{P}^\pi)^k \mathbf{R}^\pi.
\label{eq:neumann}
\end{equation}
\end{theorem}

\begin{proof}
Since $\mathbf{P}^\pi$ is row-stochastic, $\|\mathbf{P}^\pi\|_\infty=1$,
hence $\|\gamma\mathbf{P}^\pi\|_\infty=\gamma<1$. The Neumann series
$\sum_{k\ge0}\gamma^k(\mathbf{P}^\pi)^k$ therefore converges absolutely in
the induced $\infty$-norm and equals $(\mathbf{I}-\gamma\mathbf{P}^\pi)^{-1}$.
Applying it to \eqref{eq:bellmanexp} gives the claim.
\end{proof}

\begin{remark}
The Neumann series has a reading we will need repeatedly:
$(\mathbf{I}-\gamma\mathbf{P}^\pi)^{-1}$ is, up to normalisation by
$(1-\gamma)$, the matrix of \emph{discounted state occupancies}. Value is
an occupancy-weighted sum of rewards. That single observation is the
bridge from classical IRL to GAIL, and we develop it properly in
Section~\ref{sec:occupancy}.
\end{remark}

\subsection{Optimality is a fixed point}

\begin{definition}[Bellman optimality operator]
$(\mathcal{B}^\star V)(s) \triangleq \max_{a\in\A}\big[R(s,a) +
\gamma\sum_{s'}P(s'\mid s,a)V(s')\big]$.
\end{definition}

\begin{theorem}[Contraction]\label{thm:contraction}
$\|\mathcal{B}^\star V_1 - \mathcal{B}^\star V_2\|_\infty \le \gamma \|V_1 -
V_2\|_\infty$. Consequently $\mathcal{B}^\star$ has a unique fixed point
$V^\star$, value iteration converges geometrically,
$\|V^{(k)}-V^\star\|_\infty \le \gamma^k \|V^{(0)}-V^\star\|_\infty$, and
any greedy policy with respect to $V^\star$ is optimal.
\end{theorem}

\begin{proof}
For each $s$, $|\max_a f(a) - \max_a g(a)| \le \max_a |f(a)-g(a)|$; apply
this with $f(a)-g(a) = \gamma\sum_{s'}P(s'\mid s,a)(V_1-V_2)(s')$ and bound
the sum by $\gamma\|V_1-V_2\|_\infty$ using $\sum_{s'}P=1$. Banach's
fixed-point theorem gives existence, uniqueness and the rate.
\end{proof}

\subsection{Why the forward direction is easy, in four bullets}

\begin{graybox}{The Four Pillars of Forward RL}
\begin{enumerate}
\item \textbf{The unknown is finite-dimensional and the equations are
exact.} $V^\star \in \R^{|\S|}$ solves a finite system.
\item \textbf{The solution operator is a contraction.} Uniqueness and a
geometric convergence rate come for free.
\item \textbf{Expectations are finite sums over a known kernel.} Every
$\sum_{s'}P(s'\mid s,a)(\cdot)$ can be evaluated exactly.
\item \textbf{The objective is given.} There is a ground truth against
which any answer can be scored.
\end{enumerate}
\end{graybox}

Pillar 4 is what IRL removes, and that is Coordinate 1. Pillars 1--3 are
what large-scale or unknown environments remove, and that is Coordinate 2.
It is worth noticing that forward RL already has a well-developed answer to
losing pillars 1--3, which we now recall in a single page, because IRL
inherits it wholesale.

\subsection{Model-based and model-free forward RL}

\begin{enumerate}
\item \textbf{Model-based.} $P$ and $R$ are known. Solve by
\eqref{eq:neumann}, or by value/policy iteration. Value iteration converges
linearly with factor $\gamma$; policy iteration terminates finitely (at
most $|\A|^{|\S|}$ policies) and is typically far faster in practice.
\item \textbf{Model-free.} $P$ is unknown; the agent may only sample
transitions. Temporal-difference methods ($Q$-learning, SARSA) replace the
exact expectation $\sum_{s'}P(\cdot)$ by sampled successors. Tabular
$Q$-learning converges almost surely to $Q^\star$ under the Robbins--Monro
conditions $\sum_t \alpha_t=\infty$, $\sum_t \alpha_t^2<\infty$ with every
$(s,a)$ visited infinitely often \cite{watkins1992,tsitsiklis1994}. With
function approximation---policy-gradient methods included
\cite{suttonbarto2018}---the combination of bootstrapping, off-policy
sampling and approximation---the \emph{deadly triad}---can destroy the
contraction, and neither uniqueness nor convergence is guaranteed.
\end{enumerate}

\begin{remark}[Ill-conditioning as $\gamma\to1$]
In the $\infty$-norm, $\kappa_\infty(\mathbf{I}-\gamma\mathbf{P}^\pi) \le
\frac{1+\gamma}{1-\gamma}$, so the linear solve degrades as $\gamma\to
1^-$. The same $(1-\gamma)^{-1}$ factors reappear in sample-complexity
bounds: obtaining an $\varepsilon$-optimal policy with a generative model
costs $\widetilde{\mathcal{O}}\big(\frac{|\S||\A|}{\varepsilon^2(1-\gamma)^3}\big)$
samples \cite{azar2013}. The effective horizon $(1-\gamma)^{-1}$ is the
quantity that hurts, in every regime, forward and inverse alike.
\end{remark}

\clearpage
% =====================================================================
\part{The Ill-Posed Direction}\label{part:illposed}
% =====================================================================

\section{Setting Up the Inverse Problem}

In many domains---driving, surgery, aerobatic flight, dialogue---writing
down a reward is harder than demonstrating the behaviour. Hand-specified
rewards are routinely mis-specified, and optimising a mis-specified reward
produces the failure mode known as reward hacking. IRL proposes to infer
the objective from behaviour instead.

\begin{graybox}{The Two Directions}
\begin{center}
\small\setlength{\tabcolsep}{4pt}%
\begin{tabular}{@{}lccc@{}}
\textbf{Forward RL:} & dynamics $\mathcal{M}$ + reward $R$ &
$\xrightarrow{\text{optimise}}$ & optimal policy $\pi^\star$\\[4pt]
\textbf{Inverse RL:} & dynamics $\mathcal{M}$ + demonstrations $\mathcal{D}_E$ &
$\xrightarrow{\text{infer}}$ & reward $R$ (then, usually, $\pi$)
\end{tabular}
\end{center}
\end{graybox}

We observe $m$ trajectories $\mathcal{D}_E =
\{\tau^{(1)},\ldots,\tau^{(m)}\}$, with $\tau=(s_0,a_0,s_1,a_1,\ldots,s_T)$,
and we define the discounted trajectory return $R(\tau) \triangleq
\sum_{t\ge0}\gamma^t R(s_t,a_t)$.

\begin{assumption}[Expert (near-)optimality]\label{ass:expert}
$\mathcal{D}_E$ was generated by a policy $\pi_E$ that is optimal, or
near-optimal, for some unknown reward $R^\star$ in the given dynamics
$\mathcal{M}$.
\end{assumption}

Assumption~\ref{ass:expert} is the \emph{behaviour model} in its crudest
form. Taken literally it yields only inequalities:
\begin{equation}
J_R(\pi_E) \;\ge\; J_R(\pi) \qquad \text{for all } \pi\in\Pi.
\label{eq:crudeineq}
\end{equation}
A system of inequalities is a weak thing to estimate from. Much of the
technical content of modern IRL consists of replacing \eqref{eq:crudeineq}
by a \emph{probabilistic} behaviour model, which turns inequalities into a
likelihood and therefore into something a gradient can act on. We will see
exactly that happen in Section~\ref{sec:maxent}.

\section{What Exactly Can Be Inverted?}
\label{sec:invertible}

Calling IRL ``inverse optimisation'' invites the wrong intuition, because
ordinary inversion presumes injectivity. The forward map
\begin{equation}
R \;\longmapsto\; \Pi^\star(R) \;=\; \{\text{policies optimal for } R\}
\label{eq:forwardmap}
\end{equation}
is drastically many-to-one. Before choosing an algorithm we should know
precisely what information the map destroys.

\begin{graybox}{Three Objects That Must Not Be Confused}
\begin{enumerate}
\item $V^\star$ \emph{is} unique, given $\mathcal{M}$ and $R$: it is the
fixed point of a contraction.
\item $\pi^\star$ need \emph{not} be unique: ties in the greedy step give a
set $\Pi^\star(R)$.
\item $R$ is \emph{far} from unique, given behaviour: $\Pi^\star$ has
enormous fibres.
\end{enumerate}
Uniqueness of the forward solution says nothing whatever about uniqueness
of the inverse. These three statements are about three different maps.
\end{graybox}

\subsection{The solution set is a polyhedral cone}

Ng and Russell gave the exact characterisation for a finite MDP and a
fixed deterministic policy \cite{ngrussell2000}.

\begin{theorem}[Characterisation of consistent rewards, Ng \& Russell
2000]\label{thm:ngrussell}
Let $\S$ be finite, let $R\in\R^{|\S|}$ be a state-only reward, and let
$\mathbf{P}_a$ denote the matrix with entries $P(s'\mid s,a)$. A
deterministic policy $\pi$, which we may relabel so that $\pi(s)\equiv
a_1$, is optimal for $R$ if and only if
\begin{equation}
(\mathbf{P}_{a_1}-\mathbf{P}_a)(\mathbf{I}-\gamma\mathbf{P}_{a_1})^{-1}R \;\succeq\; 0
\qquad \text{for all } a\in\A\setminus\{a_1\}.
\label{eq:ngrussell}
\end{equation}
\end{theorem}

Two consequences are immediate and worth dwelling on.

\begin{corollary}\label{cor:cone}
The set of rewards consistent with an observed policy is a closed
\emph{convex cone}: it is closed under non-negative scaling and under
addition. In particular it always contains $R\equiv0$, for which every
policy is optimal, so observing behaviour alone can never exclude the
degenerate explanation ``nothing matters''.
\end{corollary}

\begin{proof}[Sketch of the degeneracies]
Beyond $R\equiv0$: any constant $R\equiv c$ makes all policies optimal;
and for any potential $\Psi:\S\to\R$, the shaped reward
\begin{equation}
R'(s,a,s') = R(s,a) + \gamma\Psi(s') - \Psi(s)
\label{eq:shaping}
\end{equation}
satisfies $J_{R'}(\pi) = J_R(\pi) - \E_{s_0\sim d_0}[\Psi(s_0)]$ for every
$\pi$, so it shifts all policy values by the same constant and preserves
the entire preference ordering over policies---hence $\Pi^\star(R') =
\Pi^\star(R)$.
\end{proof}

\begin{remark}[Correct attribution]
The optimality-preserving property of potential-based shaping
\eqref{eq:shaping} is due to Ng, Harada and Russell \cite{ngharada1999},
not the 2000 IRL paper; that work also proves the converse, that
potential-based shaping is the \emph{only} form of additive reward
modification that preserves optimality for \emph{all} transition kernels.
This converse is exactly why potential shaping is the right notion of
``nuisance parameter'' in IRL, and it is the ambiguity that AIRL sets out
to quotient away in Section~\ref{sec:airl}.
\end{remark}

\subsection{A minimal example}

Consider a single state and a single decision between \texttt{left} and
\texttt{right}; the expert chooses \texttt{left}. All we learn is
\begin{equation}
R(s,\texttt{left}) \;\ge\; R(s,\texttt{right}).
\label{eq:minimal}
\end{equation}
The pairs $(1,0)$, $(10,-3)$ and $(0,-100)$ are equally consistent. One
observation cannot pin down two numbers, and no amount of repetition of
the same observation will help: additional demonstrations add
inequalities, which sharpen the \emph{preference} structure but never fix
an absolute scale. Notice also that this example is already statistically
informative in one respect---it does rule out $R(s,\texttt{left}) <
R(s,\texttt{right})$. IRL is not hopeless; it is \emph{partially
identifying}.

\subsection{The honest target}

Given all this, the scientifically defensible outputs of IRL are:
\begin{enumerate}
\item an \textbf{equivalence class} of rewards inducing the observed
behaviour;
\item a \textbf{single representative} chosen by an explicit selection
principle (maximum margin, maximum entropy, a Bayesian prior, sparsity);
\item or simply a \textbf{policy} whose occupancy measure matches the
expert's, with no claim about reward at all.
\end{enumerate}

Item 3 is imitation learning, and it is often what one actually wants.
Item 2 is what most IRL papers deliver. Item 1 is what the data actually
supports. Confusing 2 with 1 is the single most common overclaim in the
literature; recent work makes the partial-identifiability picture precise
and shows that even with infinite data and a known model, the reward is
only determined up to explicitly characterisable transformations
\cite{cao2021,skalse2023}.

\begin{redbox}{A claim to be suspicious of}
\emph{``Our method recovers the true reward function.''} No method
recovers a reward function from behaviour alone. Methods recover a
representative of an equivalence class, selected by an assumption the
paper made. The useful question is never ``is the reward right?'' but ``is
the reward right \emph{for the use I have in mind}?'' --- and the two most
common uses, reproducing behaviour in the same MDP and transferring
behaviour to a different MDP, need very different amounts of
identification. Section~\ref{sec:airl} is entirely about the gap between
them.
\end{redbox}

\subsection{Choosing a reward class}

The unknown $R:\S\times\A\to\R$ is a function; if $\S\times\A$ is
continuous it is an infinite-dimensional object. A finite vector $w\in
\R^d$ appears only after we \emph{choose} a model class, most commonly the
linear one
\begin{equation}
R_w(s,a) = w^\top \phi(s,a).
\label{eq:linreward}
\end{equation}
This is a modelling commitment on Coordinate 1, not a law of IRL.
Alternatives include kernel and Gaussian-process rewards, neural rewards
$R_\theta$, and reward classes defined implicitly by a discriminator
architecture. We take up the question of \emph{whether $\phi$ must be
hand-designed} in full in Section~\ref{sec:exitb}; for now we simply
record what \eqref{eq:linreward} buys, because the next two sections lean
on it entirely.

\begin{proposition}[Linearity collapses return into a finite statistic]
Under \eqref{eq:linreward},
\begin{equation}
J_{R_w}(\pi) = \E_\pi\Big[\sum_{t\ge0}\gamma^t w^\top\phi(s_t,a_t)\Big]
= w^\top \E_\pi\Big[\sum_{t\ge0}\gamma^t \phi(s_t,a_t)\Big] = w^\top \mu(\pi).
\label{eq:linreturn}
\end{equation}
\end{proposition}

The vector $\mu(\pi)\in\R^d$ is a \emph{sufficient statistic} of the
policy for the purpose of evaluating any reward in the class. Everything
else about $\pi$ is irrelevant to $w$. This is the algebraic fact that
makes classical IRL tractable, and losing it is the price of leaving the
linear class.

\clearpage
% =====================================================================
\part{Two Selection Principles, Computed Exactly}\label{part:selection}
% =====================================================================

Throughout this part we stay in the \emph{safe corner}: $\S$ and $\A$
finite, the kernel $P$ fully known, and $\phi$ fixed in advance.
Coordinate 2 therefore costs us nothing---every expectation below is a
finite sum we can actually evaluate---and we can concentrate wholly on
Coordinate 1. This is the only part of the notes in which every number can
be checked by hand, and it is worth the patience: the objects defined here
reappear, in approximated form, in every method that follows.

\section{Selection Principle I: Maximum Margin}
\label{sec:selectionprinc1}

Abbeel and Ng \cite{abbeelng2004} resolve the ambiguity of
Section~\ref{sec:invertible} with a geometric rule: among all rewards
consistent with the data, prefer the one under which the expert beats the
alternatives by the largest margin.

\subsection{Feature expectations}

\begin{definition}[Feature expectation]
$\mu(\pi) \triangleq \E_{\tau\sim p_\pi}[\Phi(\tau)] = \E_\pi\big[\sum_{t\ge0}
\gamma^t \phi(s_t,a_t)\big] \in \R^d$, and the empirical expert version is
$\mu_E \triangleq \frac1m\sum_{i=1}^m \Phi(\tau^{(i)})$.
\end{definition}

\begin{lemma}[Feature matching suffices]\label{lem:featurematch}
If $\|\mu(\pi)-\mu_E\|_2 \le \varepsilon$ then for every $w$ with
$\|w\|_2\le1$,
\begin{equation}
\big|J_{R_w}(\pi) - w^\top \mu_E\big| = \big|w^\top(\mu(\pi)-\mu_E)\big| \le \varepsilon
\label{eq:featurematch}
\end{equation}
by Cauchy--Schwarz. \emph{Matching feature expectations therefore matches
expected return} simultaneously for every reward in the linear class,
\emph{without knowing which one is the truth.}
\end{lemma}

Lemma~\ref{lem:featurematch} is the reason apprenticeship learning can
sidestep identification. It also states its own limitation with unusual
clarity: the guarantee is quantified over $w\in\R^d$, and says nothing
about rewards outside $\mathrm{span}\{\phi_1,\ldots,\phi_d\}$.

\subsection{The optimisation}

Because scaling $w\mapsto\alpha w$ scales any margin, a norm constraint is
needed to make the problem bounded.

\begin{graybox}{Max-Margin IRL (Abbeel \& Ng, 2004)}
\begin{align}
\max_{\xi\in\R,\, w\in\R^d} \quad & \xi \notag\\
\text{subject to} \quad & w^\top \mu_E \;\ge\; w^\top \mu(\pi_i) + \xi, \quad i=0,\ldots,k-1, \label{eq:maxmargin}\\
& \|w\|_2 \le 1. \notag
\end{align}
\end{graybox}

This is the hard-margin SVM in disguise, with $\mu_E$ the single positive
example and $\{\mu(\pi_i)\}$ the negatives. Its dual is the
minimum-distance problem
\begin{equation}
\min_{\bar\mu \in \mathrm{conv}\{\mu(\pi_0),\ldots,\mu(\pi_{k-1})\}} \|\mu_E - \bar\mu\|_2,
\qquad \text{with optimal} \quad w \propto \mu_E - \bar\mu,
\label{eq:maxmargindual}
\end{equation}
which is the geometric content of the algorithm: the reward direction
always points from the best current imitation towards the expert. Solving
\eqref{eq:maxmargindual} directly is Abbeel and Ng's ``projection''
variant and avoids a QP solver.

\subsection{The algorithm, and what it actually returns}

\begin{algorithm}[H]
\caption{Max-Margin / Apprenticeship Learning (Abbeel \& Ng, 2004)}
\begin{algorithmic}[1]
\Require dynamics $\mathcal{M}$, features $\phi$, expert data
$\mathcal{D}_E$, tolerance $\varepsilon>0$
\State $\mu_E \gets \frac1m\sum_{i=1}^m\sum_t \gamma^t \phi(s_t^{(i)},a_t^{(i)})$
\State pick any $\pi_0$; compute $\mu(\pi_0)$ exactly by solving
$(\mathbf{I}-\gamma\mathbf{P}^{\pi_0})^{-1}$ once per feature coordinate
\For{$k=1,2,\ldots$}
  \State (\emph{reward player}) solve \eqref{eq:maxmargin} over
  $\{\mu(\pi_0),\ldots,\mu(\pi_{k-1})\}$ to get $(w^{(k)},\xi^{(k)})$
  \If{$\xi^{(k)} \le \varepsilon$} \textbf{break} \EndIf
  \State (\emph{policy player}) $\pi_k \gets$ exact forward RL solution for
  $R^{(k)} = (w^{(k)})^\top\phi$
  \State compute $\mu(\pi_k)$ exactly
\EndFor
\State solve $\lambda^\star = \argmin_{\lambda\in\Delta_k} \|\mu_E -
\sum_i \lambda_i \mu(\pi_i)\|_2$
\State \Return the \emph{mixed} policy $\tilde\pi$: sample $i\sim
\lambda^\star$ once at $t=0$, then follow $\pi_i$ forever
\end{algorithmic}
\label{alg:maxmargin}
\end{algorithm}

\begin{redbox}{The output is a mixture, not a single policy}
Algorithm~\ref{alg:maxmargin} does \emph{not} return the last $\pi_k$. Its
guarantee is for the mixed policy $\tilde\pi$, and the reason is exactly
\eqref{eq:maxmargindual}: what is provably close to $\mu_E$ is a point in
the \emph{convex hull} of the $\mu(\pi_i)$, and convex combinations of
feature expectations are realised by randomising over policies, not by any
single one of them. This works because occupancy measures mix
linearly---$\rho_{\tilde\pi} = \sum_i \lambda_i \rho_{\pi_i}$---a fact we
exploit again in Section~\ref{sec:occupancy}. Presentations that return
``the final policy'' have quietly dropped the theorem.
\end{redbox}

\begin{theorem}[Convergence, informal]
If $\phi:\S\times\A\to[0,1]^d$ then $\|\mu(\pi)\|_2 \le
\sqrt{d}/(1-\gamma)$ for all $\pi$, and Algorithm~\ref{alg:maxmargin}
produces $\|\mu(\tilde\pi)-\mu_E\|_2 \le \varepsilon$ after
$k=O\big(\frac{d}{(1-\gamma)^2\varepsilon^2}\log\frac{d}{(1-\gamma)\varepsilon}\big)$
iterations---notably, a bound independent of $|\S|$. Combined with
Lemma~\ref{lem:featurematch}, $\tilde\pi$ then performs within
$\varepsilon$ of the expert under every reward in the class.
\end{theorem}

\subsection{Limitations, stated precisely}

\begin{enumerate}
\item \textbf{Ties are broken arbitrarily.} Many $(w,\xi)$ may achieve the
same margin; the QP picks one. The selection principle is a convenience,
not a discovery.
\item \textbf{Sub-optimal experts break it.} The constraints assert the
expert is best. A noisy expert makes the feasible set shrink, or empty,
and there is no likelihood to absorb the noise. This is precisely the gap
maximum entropy fills.
\item \textbf{The inner solve is a full forward RL problem} per iteration.
In the safe corner it is exact; outside it, it is the bottleneck.
\item \textbf{It is silent about which reward is right.} Its guarantee is
about \emph{performance}, quantified over the whole class.
\end{enumerate}

\begin{navybox}{Ledger: Max-Margin IRL}
\textbf{Assumed (Coordinate 1).} Expert is exactly optimal; reward is
linear in a fixed $\phi$; ties broken by maximal margin.\\
\textbf{Approximated (Coordinate 2).} Nothing, in the safe corner: $\mu(\pi)$
comes from an exact linear solve and the inner RL is exact. Cost is one
full forward RL solve and one QP per iteration.
\end{navybox}
\begin{navyplain}
\textbf{Delivers.} A mixed policy matching $\mu_E$, plus a reward
direction $w$ whose interpretation is weak.
\end{navyplain}

\section{Selection Principle II: Maximum Entropy}
\label{sec:maxent}

Max-margin IRL treats the expert as infallible. Real demonstrations are
noisy, and worse, \emph{systematically} noisy: an expert who is
indifferent between two routes will take each about half the time, and a
method built on strict inequalities has nothing to say about that. Ziebart
et al.\ \cite{ziebart2008} replace the inequality model by a probabilistic
one, and then select among the surviving models by maximum entropy, in the
spirit of Jaynes's original principle \cite{jaynes1957}.

\begin{graybox}{The setting for this section}
Throughout Section~\ref{sec:maxent} we take a \textbf{finite} $\S$ and
$\A$, a \textbf{finite horizon} $T$, a \textbf{known} kernel $P$, and---following
the original paper---we set $\gamma=1$, writing $\Phi(\tau) =
\sum_{t=0}^{T-1}\phi(s_t,a_t)$. Discounting and infinite horizons are
restored in Section~\ref{sec:exita}, where we explain why they are
genuinely more delicate than simply inserting $\gamma^t$.
\end{graybox}

\subsection{Why a distribution over trajectories, and over what?}

We want a model $p(\tau)$ that (i) reproduces the expert's feature
statistics and (ii) commits to nothing else. The second requirement is
what ``maximum entropy'' formalises. But entropy of \emph{what}, exactly?
This is where most treatments go wrong, so we go slowly.

A trajectory is produced by two agencies: the learner, who chooses
actions, and the environment, who chooses successor states. If we maximise
plain Shannon entropy, $-\sum_\tau p(\tau)\log p(\tau)$, we implicitly
declare every trajectory equally plausible \emph{a priori}, including
trajectories the environment would essentially never produce. That is
obviously wrong. The fix is to measure entropy relative to a reference
measure encoding what the environment does.

\begin{definition}[Dynamics-induced reference measure]\label{def:refmeasure}
For $\tau = (s_0,a_0,\ldots,s_T)$ define
\begin{equation}
q(\tau) \triangleq d_0(s_0) \prod_{t=0}^{T-1} P(s_{t+1}\mid s_t,a_t).
\label{eq:refmeasure}
\end{equation}
\end{definition}

\begin{remark}
$q$ is a \emph{measure}, not a probability distribution: summing over
$\T$ gives $\sum_\tau q(\tau) = |\A|^T$, because $q$ integrates out the
environment but not the actions. It is exactly the ``uniform over your own
choices, honest about nature's'' prior. Nothing below requires $q$ to be
normalised.
\end{remark}

\begin{definition}[Relative entropy objective]
\begin{equation}
\mathcal{H}_q(p) \triangleq -\sum_{\tau\in\T} p(\tau)\log\frac{p(\tau)}{q(\tau)}
\;=\; -\KL(p\,\|\,q).
\label{eq:relentropy}
\end{equation}
\end{definition}

When $P$ is deterministic and $d_0$ uniform, $q$ is constant on feasible
trajectories and $\mathcal{H}_q$ reduces to Shannon entropy up to an
additive constant. That special case is the one usually presented; it is
not the general case.

\begin{redbox}{Two errors this definition prevents}
\textbf{(1) The retrofit error.} If one maximises plain Shannon entropy,
the stationary solution is $p(\tau)\propto e^{w^\top\Phi(\tau)}$ with
\emph{no dynamics factor}. Many expositions derive that, and then later
assert without comment that ``really'' $p(\tau)\propto
q(\tau)e^{w^\top\Phi(\tau)}$. These are solutions to two different
programs. Putting $q$ in the objective from line one makes the derivation
and the final formula agree.

\textbf{(2) The continuum error.} If one starts with $-\int p\log p\,d\tau$
over a continuous trajectory space, the objective is a \emph{differential}
entropy: it is not invariant under reparameterisation of $\tau$ and it is
not bounded above. The relative entropy $-\KL(p\|q)$ has neither defect. A
further reason to begin discrete is that the pathology is invisible until
you try to make it rigorous.
\end{redbox}

\subsection{The primal program}

\begin{equation}
\begin{aligned}
\max_p \quad & \mathcal{H}_q(p) = -\sum_\tau p(\tau)\log\frac{p(\tau)}{q(\tau)}\\
\text{s.t.} \quad & \sum_\tau p(\tau)\Phi(\tau) = \mu_E \qquad \text{(feature matching)}\\
& \sum_\tau p(\tau) = 1, \qquad p(\tau) \ge 0 \qquad \text{(normalisation)}
\end{aligned}
\label{eq:maxentprimal}
\end{equation}

The feasible set is the intersection of the simplex with $d$ affine
constraints, and $\mathcal{H}_q$ is strictly concave on it, so the
solution---if the feasible set is non-empty---is unique. Note this
uniqueness is over \emph{distributions}, not over rewards; Section~\ref{sec:invertible}
still applies.

\subsection{Variational derivation}

\textbf{Step 1: the Lagrangian.} Introduce $w\in\R^d$ for feature matching
and $c\in\R$ for normalisation:
\begin{equation}
\mathcal{L}(p,w,c) = -\sum_\tau p(\tau)\log\frac{p(\tau)}{q(\tau)}
+ w^\top\Big(\sum_\tau p(\tau)\Phi(\tau) - \mu_E\Big)
+ c\Big(\sum_\tau p(\tau) - 1\Big).
\label{eq:maxentlagrangian}
\end{equation}

\textbf{Step 2: stationarity.} Differentiating with respect to the scalar
$p(\tau)$ for a fixed $\tau$,
\begin{equation}
\frac{\partial \mathcal{L}}{\partial p(\tau)} = -\log\frac{p(\tau)}{q(\tau)}
- 1 + w^\top\Phi(\tau) + c = 0.
\label{eq:maxentstationarity}
\end{equation}

\textbf{Step 3: solve.}
\begin{equation}
p(\tau) = q(\tau)\,\exp\big(w^\top\Phi(\tau) + c - 1\big).
\label{eq:maxentsolve}
\end{equation}

\textbf{Step 4: normalise.} Imposing $\sum_\tau p(\tau)=1$ gives
$e^{c-1} = 1/Z(w)$ with

\begin{definition}[Partition function]
\begin{equation}
Z(w) \triangleq \sum_{\tau\in\T} q(\tau)\,e^{w^\top\Phi(\tau)}.
\label{eq:partitionfn}
\end{equation}
\end{definition}

\begin{graybox}{The maximum-entropy trajectory model}
\begin{equation}
p_w(\tau) = \frac{1}{Z(w)} \underbrace{q(\tau)}_{\text{what nature does}}
\;\underbrace{e^{w^\top\Phi(\tau)}}_{\text{what the agent wants}}, \qquad
Z(w) = \sum_\tau q(\tau)e^{w^\top\Phi(\tau)}.
\label{eq:maxentmodel}
\end{equation}
Two trajectories with equal reward need \emph{not} be equally likely: the
one the environment resists is rarer. Exponentially better trajectories
are exponentially more likely, but never certain---which is precisely the
graceful handling of a sub-optimal expert that max-margin lacked.
\end{graybox}

\subsection{The dual is maximum likelihood, and it is convex}

Substituting \eqref{eq:maxentmodel} back into \eqref{eq:maxentlagrangian}:
since $\log p_w(\tau) - \log q(\tau) = w^\top\Phi(\tau) - \log Z(w)$, the
entropy term becomes $-w^\top\sum_\tau p\Phi + \log Z(w)$, the
normalisation term vanishes, and the feature term cancels the first
piece, leaving
\begin{equation}
g(w) \triangleq \log Z(w) - w^\top \mu_E, \qquad \text{to be
\emph{minimised} over } w.
\label{eq:maxentdual}
\end{equation}

\begin{theorem}[Duality and maximum likelihood]\label{thm:maxentduality}
Minimising \eqref{eq:maxentdual} is exactly maximising the average
log-likelihood of the demonstrations under $p_w$:
\begin{equation}
\begin{aligned}
\frac1m\sum_{i=1}^m \log p_w(\tau^{(i)})
&= \frac1m\sum_{i=1}^m \Big[\log q(\tau^{(i)}) + w^\top\Phi(\tau^{(i)}) - \log Z(w)\Big]\\
&= \underbrace{w^\top\mu_E - \log Z(w)}_{=-g(w)}
+ \underbrace{\tfrac1m\sum_i \log q(\tau^{(i)})}_{\text{does not involve } w}.
\end{aligned}
\label{eq:maxentmle}
\end{equation}
\end{theorem}

The dropped term is the log-likelihood of the observed transitions under
the true dynamics; it is a constant of the optimisation, which is why it
is usually invisible.

\begin{theorem}[Convexity]\label{thm:logzconvex}
$w\mapsto \log Z(w)$ is convex, with
\begin{equation}
\nabla_w \log Z(w) = \E_{p_w}[\Phi(\tau)], \qquad
\nabla_w^2 \log Z(w) = \mathrm{Cov}_{p_w}[\Phi(\tau)] \succeq 0.
\label{eq:logzgrad}
\end{equation}
Hence $g$ is convex; it is strictly convex, and the optimal $w$ unique,
precisely when the features $\Phi$ are not almost-surely confined to a
proper affine subspace.
\end{theorem}

\begin{proof}
$\nabla_w Z = \sum_\tau q(\tau)e^{w^\top\Phi}\Phi$, so $\nabla_w \log Z =
\sum_\tau p_w(\tau)\Phi(\tau)$. Differentiating once more gives
$\E_{p_w}[\Phi\Phi^\top] - \E_{p_w}[\Phi]\E_{p_w}[\Phi]^\top$, a covariance
matrix.
\end{proof}

\begin{corollary}[The learning rule]
\begin{equation}
\boxed{-\nabla_w g(w) \;=\; \mu_E - \E_{p_w}[\Phi(\tau)]}
\label{eq:learningrule}
\end{equation}
\end{corollary}

Gradient ascent moves $w$ so as to close the gap between what the expert
did and what the current model predicts. At the optimum the two feature
expectations coincide.

\begin{remark}[A pleasing tension worth noticing]
Theorem~\ref{thm:logzconvex} says the reward \emph{parameter} $w$ has a
unique optimum. Section~\ref{sec:invertible} says the reward
\emph{function} is not identified. Both are true, and the resolution is
that maximum entropy is a selection principle: it picks one point out of
the equivalence class, and picks it uniquely. Uniqueness of an estimator
is not identifiability of an estimand.
\end{remark}

\subsection{Computing $Z$: the whole difficulty, made concrete}
\label{sec:computingz}

Equation~\eqref{eq:learningrule} looks like a two-line algorithm, and the
entire remainder of these notes exists because of the innocuous term
$\E_{p_w}[\Phi]$. Evaluating it requires $Z(w)$, and $Z(w)$ is a sum over
$\T$.

\textbf{Enumeration is hopeless but finite.} With $|\A|$ actions and
horizon $T$, $|\T|$ grows like $|\A|^T$ (times the branching of $P$). For
a $5\times5$ gridworld with 4 actions and $T=20$ that is $4^{20}\approx
10^{12}$ trajectories. Finite, and still impossible. The value of the
discrete formulation is not that we enumerate; it is that $Z$ is
unambiguously \emph{defined}, so we can look for structure.

\textbf{The structure is that $q$ factorises.} Because $q(\tau) =
d_0(s_0)\prod_t P(s_{t+1}\mid s_t,a_t)$ and $\Phi(\tau) =
\sum_t\phi(s_t,a_t)$, the sum over trajectories telescopes into a backward
recursion over states. Write $R=R_w$ for brevity.

\begin{proposition}[Exact backward recursion for $Z$]\label{prop:backwardz}
Define $Z_T(s) \triangleq 1$ and, for $t=T-1,\ldots,0$,
\begin{equation}
Z_t(s) \;=\; \sum_{a\in\A} e^{R(s,a)} \sum_{s'\in\S} P(s'\mid s,a)\,Z_{t+1}(s').
\label{eq:backwardz}
\end{equation}
Then $Z(w) = \sum_s d_0(s) Z_0(s)$. Writing $V_t(s) = \log Z_t(s)$, this is
\begin{equation}
V_t(s) = \log\sum_a \exp\Big(R(s,a) + \underbrace{\log\sum_{s'}P(s'\mid
s,a)e^{V_{t+1}(s')}}_{=:\,Q_t(s,a)}\Big).
\label{eq:backwardv}
\end{equation}
The cost is $O(T|\S|^2|\A|)$ instead of $O(|\A|^{|\S|})$.
\end{proposition}

\begin{proof}
Let $Z_t(s)$ be the sum, over all suffixes $(a_t,s_{t+1},\ldots,s_T)$
started at $s_t=s$, of $\prod_{u\ge t} P(s_{u+1}\mid s_u,a_u)
\exp\big(\sum_{u\ge t} R(s_u,a_u)\big)$. Splitting off the first action and
successor gives exactly \eqref{eq:backwardz}; the base case is the empty
suffix. Averaging over $s_0\sim d_0$ gives $Z(w)$.
\end{proof}

\begin{graybox}{Answer: where does knowledge of $P$ enter?}
Look at \eqref{eq:backwardz}. The kernel $P(s'\mid s,a)$ is consumed
inside the inner sum, once per $(s,a)$ per timestep. This is the
\textbf{first} of exactly two places where MaxEnt IRL needs a model. The
second appears in Section~\ref{sec:forwardpass}. When we ask, in
Section~\ref{sec:exitc}, ``what if $P$ is unknown?'', these two sums are
the entire content of the question.
\end{graybox}

\subsection{A worked example you can check by hand}

\begin{example}[Two states, two actions, horizon two]\label{ex:worked}
Let $\S=\{s_1,s_2\}$, $\A=\{L,R\}$, $d_0=\delta_{s_1}$, $T=2$, and
deterministic dynamics: $L$ moves to $s_1$, $R$ moves to $s_2$, from
either state. Let $R(s_1,L)=0$, $R(s_1,R)=1$, $R(s_2,L)=0$, $R(s_2,R)=2$.

\textbf{By enumeration.} Four action sequences:
\begin{center}
\begin{tabular}{@{}cccc@{}}
\toprule
sequence & path & return & weight $e^{\text{return}}$\\
\midrule
$LL$ & $s_1\to s_1\to s_1$ & $0+0=0$ & $1.0000$\\
$LR$ & $s_1\to s_1\to s_2$ & $0+1=1$ & $2.7183$\\
$RL$ & $s_1\to s_2\to s_1$ & $1+0=1$ & $2.7183$\\
$RR$ & $s_1\to s_2\to s_2$ & $1+2=3$ & $20.0855$\\
\midrule
& & $Z=$ & $26.5221$\\
\bottomrule
\end{tabular}
\end{center}

\textbf{By the backward recursion~\eqref{eq:backwardv}.} With $V_2\equiv0$:
\begin{align*}
V_1(s_1) &= \log(e^0+e^1) = 1.3133, & V_1(s_2) &= \log(e^0+e^2) = 2.1269,\\
Q_0(s_1,L) &= 0 + V_1(s_1) = 1.3133, & Q_0(s_1,R) &= 1 + V_1(s_2) = 3.1269,\\
V_0(s_1) &= \log(e^{1.3133}+e^{3.1269}) = 3.2780, & Z &= e^{V_0(s_1)} = 26.522. \;\checkmark
\end{align*}
And the model's first action probability agrees both ways:
\begin{equation*}
p_w(a_0=R) = e^{Q_0(s_1,R) - V_0(s_1)} = \frac{22.804}{26.522} = 0.8598
= \frac{e^1+e^3}{Z} = \frac{22.804}{26.522}. \;\checkmark
\end{equation*}
Four trajectories collapsed into four $(s,a)$ evaluations. With $T=20$ the
same collapse turns the $4^{20}\approx10^{12}$ trajectories of the
$5\times5$ gridworld into $25\times4\times20=2000$ state--action--time evaluations.
\end{example}

\subsection{The stochastic-dynamics trap, and maximum \emph{causal} entropy}
\label{sec:causaltrap}

Example~\ref{ex:worked} was deterministic, and that was not laziness. Look
again at \eqref{eq:backwardv}. The inner term
\begin{equation}
\log \sum_{s'} P(s'\mid s,a)\, e^{V_{t+1}(s')}
\label{eq:innersoftmax}
\end{equation}
is a \emph{soft maximum} over successor states, not an expectation. Since
$\log\E[e^X] \ge \E[X]$ by Jensen, it systematically over-values actions
with high variance in their outcomes.

The pathology is sharper than a bias in a formula. It is a statement about
what the model believes.

\begin{proposition}[The naive model secretly rewrites the
dynamics]\label{prop:naivetrap}
Under $p_w$ of \eqref{eq:maxentmodel}, the conditional law of the
successor state is
\begin{equation}
p_w(s_{t+1}=s' \mid s_t=s, a_t=a) = \frac{P(s'\mid s,a)\,e^{V_{t+1}(s')}}
{\sum_{s''}P(s''\mid s,a)e^{V_{t+1}(s'')}} \;\ne\; P(s'\mid s,a) \quad
\text{in general.}
\label{eq:naivetrap}
\end{equation}
\end{proposition}

\begin{proof}
Marginalise \eqref{eq:maxentmodel} over all suffixes after $s_{t+1}$ using
Proposition~\ref{prop:backwardz}; the suffix sum contributes the factor
$Z_{t+1}(s') = e^{V_{t+1}(s')}$.
\end{proof}

\begin{redbox}{What this means}
The maximum-entropy trajectory model with stochastic dynamics is not a
model of an agent choosing actions in the given MDP. It is a model of an
agent who \emph{also} tilts nature towards fortunate outcomes. Fitted to a
real expert, it will explain lucky outcomes as intentional and infer
rewards that reward luck. In a deterministic MDP the tilt is vacuous ($P$
is a point mass), which is why the original presentation is correct as
stated and why the problem is so easy to miss.
\end{redbox}

The remedy is to stop optimising over arbitrary trajectory distributions
and optimise over \emph{policies}, forcing the true dynamics to appear as
an unchangeable factor.

\begin{definition}[Causal entropy]
For a (possibly time-dependent) policy $\pi$, let $p_\pi(\tau) =
d_0(s_0)\prod_t \pi_t(a_t\mid s_t)P(s_{t+1}\mid s_t,a_t)$ and
\begin{equation}
\mathcal{H}^{\mathrm{causal}}(\pi) \triangleq \E_{p_\pi}\Big[\sum_{t=0}^{T-1}
-\log \pi_t(a_t\mid s_t)\Big].
\label{eq:causalentropy}
\end{equation}
\end{definition}

Only the agent's own choices are counted as uncertainty; randomness
supplied by nature is not credited to the agent.

\begin{theorem}[Maximum causal entropy IRL, Ziebart
2010~\cite{ziebart2010thesis}]\label{thm:maxcausal}
The solution of
\begin{equation}
\max_\pi \; \mathcal{H}^{\mathrm{causal}}(\pi) \quad \text{s.t.}\quad
\E_{p_\pi}[\Phi(\tau)] = \mu_E
\label{eq:maxcausalprogram}
\end{equation}
is the soft-optimal policy for $R_w=w^\top\phi$ with $w$ the multiplier,
given by the backward recursion
\begin{align}
Q_t(s,a) &= R_w(s,a) + \sum_{s'} P(s'\mid s,a) V_{t+1}(s'), \label{eq:causalQ}\\
V_t(s) &= \log\sum_a \exp Q_t(s,a), \label{eq:causalV}\\
\pi_t(a\mid s) &= \exp\big(Q_t(s,a) - V_t(s)\big), \label{eq:causalpi}
\end{align}
with $V_T\equiv0$. The dual objective remains convex and its gradient
remains $\mu_E - \E_{p_\pi}[\Phi]$.
\end{theorem}

\begin{graybox}{The one-symbol difference}
Compare \eqref{eq:backwardv} with \eqref{eq:causalQ}--\eqref{eq:causalV}.
Naive MaxEnt takes a \emph{soft} maximum over both actions and successor
states; maximum causal entropy takes a soft maximum over actions and a
\emph{plain expectation} over successor states:
\begin{center}
\begin{tabular}{@{}ll@{}}
naive: & $\log\sum_a \exp\big(R + \log\sum_{s'} P e^V\big)$\\
causal: & $\log\sum_a \exp\big(R + \sum_{s'} P V\big)$
\end{tabular}
\end{center}
They coincide exactly when $P$ is deterministic. Every modern method---soft
$Q$-learning, soft actor--critic, GAIL's entropy term, AIRL---uses the
causal version. When a paper says ``MaxEnt IRL,'' it almost always means
this.
\end{graybox}

\subsection{The forward pass: from a policy to expected features}
\label{sec:forwardpass}

The backward pass produced a policy. We still need $\E_{p_\pi}[\Phi]$,
which requires knowing \emph{where} that policy goes. That is a forward
sweep.

\begin{proposition}[Forward state-visitation recursion]
Let $d_t(s) = \Prb_{p_\pi}(s_t=s)$. Then $d_0$ is the initial distribution
and
\begin{equation}
d_{t+1}(s') = \sum_{s\in\S}\sum_{a\in\A} d_t(s)\,\pi_t(a\mid s)\,P(s'\mid s,a),
\label{eq:forwardrecursion}
\end{equation}
whence
\begin{equation}
\E_{p_\pi}[\Phi] = \sum_{t=0}^{T-1}\sum_{s,a} d_t(s)\,\pi_t(a\mid s)\,\phi(s,a).
\label{eq:forwardexpect}
\end{equation}
\end{proposition}

\begin{graybox}{Answer: where does knowledge of $P$ enter? (second and
last place)}
Equation~\eqref{eq:forwardrecursion}. Together with \eqref{eq:causalQ},
these are the \emph{only} two consumers of $P$ in exact MaxEnt IRL. A
backward pass computes \emph{what the agent would do}; a forward pass
computes \emph{where doing it takes them}. Both need the kernel; neither
needs anything else.
\end{graybox}

\subsection{Occupancy measures: the same object, better named}
\label{sec:occupancy}

Equation~\eqref{eq:forwardexpect} is a weighted average of $\phi$ over
state--action pairs. Naming that weighting will pay for itself many times
over.

\begin{definition}[Discounted occupancy measure]
\begin{equation}
\rho_\pi(s,a) \triangleq (1-\gamma)\sum_{t=0}^\infty \gamma^t \Prb_\pi(s_t=s,a_t=a),
\label{eq:occupancydef}
\end{equation}
a probability distribution over $\S\times\A$. Then
\begin{equation}
\mu(\pi) = \frac1{1-\gamma}\sum_{s,a}\rho_\pi(s,a)\phi(s,a), \qquad
J_R(\pi) = \frac1{1-\gamma}\sum_{s,a}\rho_\pi(s,a)R(s,a).
\label{eq:occupancylinear}
\end{equation}
\end{definition}

\begin{remark}[Occupancy is not action frequency]
$\rho_\pi$ is a distribution over \emph{pairs}. ``Brake'' in front of a
pedestrian and ``brake'' on an empty motorway are different events because
their states differ. MaxEnt IRL never merely counts action labels.
\end{remark}

Equation~\eqref{eq:occupancylinear} says return is a \emph{linear
functional} of occupancy. We will build all of Part~\ref{part:adversarial}
on that sentence.

\subsection{The complete exact algorithm}

Because the objective is convex (Theorem~\ref{thm:logzconvex}), any
diminishing-step gradient scheme converges to the global optimum in $w$.
The per-iteration cost is $O(T|\S|^2|\A|)$ for each pass---linear in the
horizon, quadratic in the number of states. That quadratic is the thing
that will not survive contact with a real problem.

\begin{algorithm}[H]
\caption{Exact Maximum Causal Entropy IRL (finite $\S,\A$, known $P$)}
\begin{algorithmic}[1]
\Require dynamics $\mathcal{M}$ with known $P$, features $\phi$, expert
data $\mathcal{D}_E$, step size $\eta$
\State $\mu_E \gets \frac1m\sum_i \Phi(\tau^{(i)})$
\State initialise $w\gets0$
\For{$k=1,2,\ldots$ until $\|\mu_E-\hat\mu\|$ small}
  \State $R_w(s,a)\gets w^\top\phi(s,a)$ for all $(s,a)$
  \State \textbf{backward pass:} $V_T\equiv0$; for $t=T-1,\ldots,0$ compute
  $Q_t,V_t,\pi_t$ by \eqref{eq:causalQ}--\eqref{eq:causalpi} \Comment{uses $P$}
  \State \textbf{forward pass:} $d_0\gets d_0$; for $t=0,\ldots,T-1$
  propagate $d_{t+1}$ by \eqref{eq:forwardrecursion} \Comment{uses $P$}
  \State $\hat\mu \gets \sum_{t,s,a} d_t(s)\pi_t(a\mid s)\phi(s,a)$
  \State $w \gets w + \eta\,(\mu_E - \hat\mu)$ \Comment{gradient~\eqref{eq:learningrule}}
\EndFor
\State \Return $w$, and the soft-optimal policy $\pi$
\end{algorithmic}
\label{alg:maxcausalent}
\end{algorithm}

\begin{navybox}{Ledger: Exact MaxEnt / MaxCausalEnt IRL}
\textbf{Assumed (Coordinate 1).} Expert is \emph{Boltzmann-rational}:
exponentially prefers better trajectories rather than always choosing the
best. Reward linear in fixed $\phi$. Among all models matching $\mu_E$,
take the maximum-causal-entropy one.\\
\textbf{Approximated (Coordinate 2).} Nothing. $Z$ is exact via a backward
pass; $\E[\Phi]$ is exact via a forward pass. Requires: finite $\S,\A$, and
the kernel $P$ in closed form.\\
\textbf{Delivers.} A unique $w$ (convex problem), a stochastic policy that
degrades gracefully with expert noise, and a likelihood---so the model can
be compared, regularised and given priors, none of which max-margin
allowed.\\
\textbf{Breaks when.} $|\S|$ is large or continuous, the horizon is
infinite, or $P$ is unknown. These are Sections~\ref{sec:exita},
\ref{sec:exitb} and \ref{sec:exitc}.
\end{navybox}

\clearpage
% =====================================================================
\part{Three Exits from the Safe Corner}\label{part:exits}
% =====================================================================

Algorithm~\ref{alg:maxcausalent} is a complete, correct, globally
convergent IRL method. It is also unusable on any problem anyone actually
cares about, for three separate reasons. Each reason is one of the
assumptions we made on entering Part~\ref{part:selection}, and each is a
question every reader of the IRL literature ends up asking. We take them
one at a time, and in each case we ask the same two things: \emph{what
survives, and what must be replaced?}

\section{Exit A: What Changes When the Problem Is Not Finite?}
\label{sec:exita}

``Not finite'' hides three separate difficulties: an unbounded horizon, an
unbounded number of states, and a continuum of actions. They fail
differently.

\subsection{Infinite horizon: discounting, and why trajectory entropy dies}

The finite-horizon derivation of Section~\ref{sec:maxent} does not survive
$T\to\infty$ as written, and the reason is not the return but the
\emph{entropy}. The causal entropy $\E\big[\sum_{t<T}-\log\pi(a_t\mid
s_t)\big]$ grows linearly in $T$ for any non-deterministic policy, so the
objective diverges. The repair is the same discounting we already apply to
reward.

\begin{definition}[Discounted causal entropy and the soft objective]
\begin{equation}
\mathcal{H}_\gamma(\pi) \triangleq \E_\pi\Big[\sum_{t=0}^\infty \gamma^t
\big(-\log\pi(a_t\mid s_t)\big)\Big], \qquad
J_R^\beta(\pi) \triangleq J_R(\pi) + \beta\,\mathcal{H}_\gamma(\pi).
\label{eq:softobjective}
\end{equation}
\end{definition}

The temperature $\beta>0$ is now explicit; Section~\ref{sec:maxent} used
$\beta=1$.

\begin{definition}[Soft Bellman operator]
\begin{equation}
(\mathcal{B}_\soft Q)(s,a) \triangleq R(s,a) + \gamma\,\E_{s'\sim P(\cdot\mid
s,a)}\Big[\beta\log\sum_{a'}\exp\big(Q(s',a')/\beta\big)\Big].
\label{eq:softbellman}
\end{equation}
\end{definition}

\begin{lemma}[Log-sum-exp is a non-expansion]\label{lem:lse}
For $x,y\in\R^{|\A|}$ and $\beta>0$, $\big|\beta\log\sum_a e^{x_a/\beta} -
\beta\log\sum_a e^{y_a/\beta}\big| \le \|x-y\|_\infty$.
\end{lemma}

\begin{proof}
Let $\delta=\|x-y\|_\infty$. Then $x_a\le y_a+\delta$ for all $a$, so
$\sum_a e^{x_a/\beta}\le e^{\delta/\beta}\sum_a e^{y_a/\beta}$; taking
$\beta\log$ gives $\beta\log\sum e^{x_a/\beta} \le \beta\log\sum
e^{y_a/\beta} + \delta$. Exchanging $x$ and $y$ gives the other
inequality.
\end{proof}

\begin{theorem}[Soft value iteration converges]
$\mathcal{B}_\soft$ is a $\gamma$-contraction in $\|\cdot\|_\infty$, so it
has a unique fixed point $Q_\soft$, soft value iteration converges
geometrically, and the soft-optimal policy is $\pi_\soft(a\mid s)\propto
\exp(Q_\soft(s,a)/\beta)$, which maximises \eqref{eq:softobjective}.
\end{theorem}

\begin{proof}
Apply Lemma~\ref{lem:lse} pointwise inside the expectation in
\eqref{eq:softbellman} and use $\sum_{s'}P=1$.
\end{proof}

\begin{remark}[The temperature interpolates]
Since
\begin{equation*}
\max_a Q(s,a) \;\le\; \beta\log\sum_a e^{Q(s,a)/\beta} \;\le\; \max_a
Q(s,a) + \beta\log|\A|,
\end{equation*}
the soft operator is within $\beta\log|\A|$ of $\mathcal{B}^\star$
everywhere. As $\beta\to0$ we recover hard optimality and max-margin's
behaviour model; as $\beta\to\infty$ the policy becomes uniform and the
demonstrations are treated as uninformative. \emph{The temperature is a
modelling choice on Coordinate 1}: it says how rational you think the
expert is. Fitting $\beta$ from data is possible but is entangled with the
scale of $R$, since $(R,\beta)$ and $(cR,c\beta)$ give the same
policy---one more instance of the identifiability theme.
\end{remark}

\begin{graybox}{What survived Exit A, part 1}
The whole structure of Section~\ref{sec:maxent} survives an infinite
horizon, provided we use \emph{discounted causal} entropy and a stationary
soft Bellman fixed point in place of a finite backward sweep. What we lose
is only the interpretation of $\log Z$ as a partition function of a finite
trajectory set; what replaces it is a soft value function. Nothing
statistical changes.
\end{graybox}

\subsection{Large or continuous state spaces}

Here the loss is real. Both passes of Algorithm~\ref{alg:maxcausalent} are
tabular sweeps costing $O(|\S|^2|\A|)$. Three things break at once.

\begin{enumerate}
\item \textbf{The tables do not fit.} $V$, $\pi$ and $d_t$ are arrays
indexed by $\S$. For continuous $\S$ they are functions and must be
approximated---by linear architectures, or in practice by neural
networks. This reintroduces every hazard of the deadly triad into the
inner loop of IRL.
\item \textbf{Discretisation does not save you.} An $\varepsilon$-net of
$[0,1]^n$ has $\varepsilon^{-n}$ cells. A ten-dimensional state at
resolution $0.1$ is $10^{10}$ cells, and the sweep is quadratic in that.
\item \textbf{Sums become integrals and the reference measure becomes
load-bearing.} In \eqref{eq:partitionfn}, $\sum_\tau q(\tau)(\cdot)$
becomes $\int(\cdot)\,dq(\tau)$. This is precisely where the
relative-entropy formulation of Definition~\ref{def:refmeasure} pays off:
$-\KL(p\|q)$ is well-defined and reparameterisation-invariant on a
continuous space, whereas $-\int p\log p$ is neither. Had we begun with
the continuous differential entropy, the objective would have changed
value under a smooth change of coordinates on $\T$---an absurdity for a
principle that is supposed to encode ignorance.
\end{enumerate}

\subsection{Continuous actions}

The soft-max over actions, $\beta\log\sum_a\exp(Q/\beta)$, becomes
$\beta\log\int_\A \exp(Q(s,a)/\beta)\,da$: an integral with no closed form
for a general $Q$. The standard resolution is variational: represent the
policy explicitly and fit it to the Boltzmann target,
\begin{equation}
\pi \;\leftarrow\; \argmin_{\pi'\in\Pi_\theta} \KL\Big(\pi'(\cdot\mid s)\,\Big\|\,
\tfrac1{Z(s)}\exp\big(Q(s,\cdot)/\beta\big)\Big),
\label{eq:sacstep}
\end{equation}
which is exactly the policy-improvement step of soft actor--critic
\cite{haarnoja2018}, and an instance of the broader control-as-inference
picture in which entropy-regularised RL is exact probabilistic inference
in a graphical model \cite{levine2018tutorial}. Note the structure: an
intractable normaliser over actions is replaced by a learned sampler. We
are about to meet the same move one level up, where an intractable
normaliser over \emph{trajectories} is replaced by a learned policy.

\subsection{Summary of Exit A}

\begin{center}
\begin{tabular}{@{}p{0.24\linewidth}p{0.34\linewidth}p{0.36\linewidth}@{}}
\toprule
\textbf{Object} & \textbf{Finite case} & \textbf{Infinite / continuous case}\\
\midrule
Entropy objective & causal entropy over $T$ steps & discounted causal
entropy; relative to $q$\\
Backward pass & exact sweep, $T$ steps & soft Bellman fixed point
(contraction survives)\\
Soft-max over actions & $\log\sum_a$ & intractable integral; variational
policy (SAC)\\
Forward pass / $\rho_\pi$ & exact array recursion & Monte Carlo rollouts or
density models\\
Partition function $Z$ & finite sum, exact via DP & intractable; estimated
or bypassed\\
Convexity of the dual in $w$ & holds & holds \emph{if} the reward stays
linear in $w$\\
\bottomrule
\end{tabular}
\end{center}

The last row deserves emphasis: leaving the safe corner along Exit A alone
does \emph{not} cost us convexity. Convexity is lost at Exit B.

\section{Exit B: Must the Features Be Hand-Crafted?}
\label{sec:exitb}

Short answer: no, and modern methods do not. But the question hides a more
useful one---\emph{what was the fixed $\phi$ doing for us?}---and the
honest answer is: quite a lot.

\subsection{What a fixed linear class buys}

\begin{enumerate}
\item \textbf{A sufficient statistic.} By Proposition~3.4, a policy is
summarised for all reward-evaluation purposes by $\mu(\pi)\in\R^d$. An
infinite-dimensional object collapses to $d$ numbers.
\item \textbf{Convexity.} Theorem~\ref{thm:logzconvex} needs $R$ linear in
the parameter. This is what makes MaxEnt IRL a well-behaved convex
estimation problem.
\item \textbf{A class-wide guarantee.} Lemma~\ref{lem:featurematch}
certifies performance for \emph{every} reward in the class at once, which
is much stronger than certifying it for one estimate.
\item \textbf{Interpretability.} ``The expert weights collision risk 8
times lane deviation'' is a sentence one can act on.
\end{enumerate}

\subsection{What it costs}

\begin{redbox}{Misspecification is invisible from inside}
Every guarantee above is quantified over $w\in\R^d$ with $\phi$ fixed. If
the true reward is not in $\mathrm{span}\{\phi_1,\ldots,\phi_d\}$,
matching $\mu_E$ implies \emph{nothing} about $J_{R^\star}$. Worse, the fit
will look fine: the residual lives in a direction the model cannot see.
Classical apprenticeship-learning results on driving worked because a
human had already encoded the task into features such as ``distance to
nearest car'' and ``off-road indicator''. The domain knowledge was in
$\phi$, not the algorithm.
\end{redbox}

\subsection{The ambiguity that appears when $\phi$ is learned}

Suppose we learn $\phi_\eta$ and $w$ jointly. For any invertible $A\in
\R^{d\times d}$,
\begin{equation}
w^\top\phi(s,a) = (A^{-\top}w)^\top\big(A\phi(s,a)\big),
\label{eq:featureambiguity}
\end{equation}
so $(w,\phi)$ and $(A^{-\top}w, A\phi)$ are indistinguishable. Only the
\emph{scalar field} $R$ and the \emph{span} of the features are
meaningful; individual coordinates of $w$ are not. Statements like ``the
model learned that feature 3 matters'' are therefore meaningless without a
convention pinning down $A$. This is a Coordinate-1 ambiguity \emph{added}
by learning features, on top of the reward ambiguity we already had.

\subsection{The good news: the learning rule does not change}

One might fear that abandoning linearity forces a new derivation. It does
not.

\begin{proposition}[MaxEnt gradient for an arbitrary reward
class]\label{prop:generalgrad}
Let $R_\theta$ be any differentiable reward parameterisation, and let
$R_\theta(\tau) = \sum_t \gamma^t R_\theta(s_t,a_t)$. With $p_\theta(\tau)
\propto q(\tau)e^{R_\theta(\tau)}$, the log-likelihood $L(\theta) =
\frac1m\sum_i R_\theta(\tau^{(i)}) - \log Z(\theta)$ has gradient
\begin{equation}
\nabla_\theta L(\theta) = \underbrace{\E_{\tau\sim\mathcal{D}_E}\big[\nabla_\theta
R_\theta(\tau)\big]}_{\text{expert}} - \underbrace{\E_{\tau\sim
p_\theta}\big[\nabla_\theta R_\theta(\tau)\big]}_{\text{current model}}.
\label{eq:generalgrad}
\end{equation}
\end{proposition}

\begin{proof}
$\nabla_\theta\log Z(\theta) = \frac1Z\sum_\tau q(\tau)e^{R_\theta(\tau)}
\nabla_\theta R_\theta(\tau) = \E_{p_\theta}[\nabla_\theta R_\theta(\tau)]$.
\end{proof}

Equation~\eqref{eq:generalgrad} is \eqref{eq:learningrule} with $\phi$
replaced by $\nabla_\theta R_\theta$. In the linear case $\nabla_w
R_w(\tau) = \Phi(\tau)$ and the two coincide exactly. \emph{Feature
matching was never the principle; matching the sufficient statistics of
the reward class was.} This observation is what makes deep MaxEnt IRL
\cite{wulfmeier2015} and guided cost learning \cite{finn2016} possible
with no new theory.

\subsection{What is lost: convexity, and the need for inductive bias}

With $R_\theta$ nonlinear, $L$ is no longer concave. We get local optima,
sensitivity to initialisation, and---most importantly---capacity to
overfit. A sufficiently expressive $R_\theta$ can memorise the
demonstrations: assign huge reward exactly on the observed state--action
pairs and zero elsewhere. That reward perfectly explains the data and is
useless.

Useful counter-pressures, all of which are Coordinate-1 commitments:

\begin{itemize}
\item \textbf{Architectural restriction.} Most powerfully, making $R$
depend on the state only, $R_\theta(s)$. We will see in
Section~\ref{sec:airl} that this is not a convenience but a necessary
condition for a transfer guarantee.
\item \textbf{Smoothness and norm penalties} on $R_\theta$; gradient
penalties on a discriminator.
\item \textbf{Bottlenecks}, including explicit information bottlenecks on
the reward representation.
\item \textbf{Multi-task structure}: several tasks sharing one
representation with different weights, which restores identifiability of
the shared part.
\item \textbf{Priors}, i.e.\ Bayesian IRL \cite{ramachandran2007}, which
makes the selection principle explicit rather than implicit.
\end{itemize}

\begin{graybox}{Answer to Exit B}
No, $\phi$ need not be hand-crafted, and the maximum-entropy learning rule
\eqref{eq:generalgrad} is unchanged when it is not. What changes is that
the estimation problem stops being convex and starts requiring
regularisation, and that a new reparameterisation ambiguity is
introduced. The choice of reward class does not disappear when features
are learned---it becomes the choice of \emph{architecture}. In
Part~\ref{part:adversarial} it becomes the choice of \emph{discriminator},
which is why we insist that a discriminator architecture is a reward
class in disguise.
\end{graybox}

\section{Exit C: What If We Do Not Know $P$?}
\label{sec:exitc}

This is the question that most often derails a first reading, because two
entirely different estimation problems are usually collapsed into one.
Separate them first.

\begin{graybox}{Two statistical jobs, frequently confused}
\textbf{Job 1 --- estimate the expert's statistics $\rho_E$ (or $\mu_E$).}
Data: the demonstrations. Needs no model and no interaction: an average
over $\mathcal{D}_E$. Its weakness is \emph{support}---it says nothing
outside the states the expert visited.

\medskip
\textbf{Job 2 --- estimate the \emph{current candidate policy's}
statistics $\rho_{\pi_w}$.}
This is not in the data at all: $\pi_w$ changes every gradient step. It
requires either a model of $P$ (to compute) or the ability to interact
with the environment (to sample). This is the job that makes IRL
expensive, and the only one that needs $P$.
\end{graybox}

Recall from Sections~\ref{sec:computingz} and \ref{sec:forwardpass} that
$P$ is consumed in exactly two places, both belonging to Job 2: the
expectation over successors in the backward pass \eqref{eq:causalQ}, and
the propagation of visitation in the forward pass
\eqref{eq:forwardrecursion}. Everything below is about replacing those two
sums.

\subsection{Four regimes}

\begin{center}
\begin{tabular}{@{}p{0.26\linewidth}p{0.42\linewidth}p{0.26\linewidth}@{}}
\toprule
\textbf{What you have} & \textbf{What to do about Job 2} & \textbf{Price}\\
\midrule
$P$ in closed form, finite $\S$ & Exact backward + forward passes
(Algorithm~\ref{alg:maxcausalent}) & $O(T|\S|^2|\A|)$ per gradient step\\
A simulator you can reset & Roll out $\pi_w$; estimate $\E[\Phi]$ by Monte
Carlo. You never need to \emph{write down} $P$ & Variance; many rollouts
per gradient step\\
Only the real environment & Fit $\widehat P$ and plan, or estimate
occupancies directly from on-policy data & Exploration cost; model bias;
safety risk\\
Only a fixed offline dataset & Job 2 is not identifiable off-support.
Requires pessimism, conservative penalties, or restricting the policy
class & Genuinely hard; an active research area\\
\bottomrule
\end{tabular}
\end{center}

\begin{remark}[A simulator is not a model, and that is the point]
The second row is the one that dissolves most of the confusion.
``Model-free'' does not mean ``dynamics-free''. It means you never need
$P$ as an object; you need only the ability to draw $s'\sim P(\cdot\mid
s,a)$. Almost every modern IRL method lives in this row.
\end{remark}

\subsection{Sampling the gradient}

The Monte Carlo replacement is direct. Draw $n$ rollouts
$\tau_1,\ldots,\tau_n$ from the current soft policy and set
\begin{equation}
\widehat{\E_{p_w}[\Phi]} = \frac1n\sum_{j=1}^n \Phi(\tau_j), \qquad
\widehat{\nabla_w L} = \mu_E - \frac1n\sum_{j=1}^n \Phi(\tau_j),
\label{eq:mcgrad}
\end{equation}
which is unbiased and consistent. The catch is that the current soft
policy must itself be obtained by an inner forward-RL solve at every outer
step---extremely wasteful, since $w$ changes only slightly.

Guided cost learning \cite{finn2016} removes the waste with importance
sampling: maintain a sampler $\pi_\samp$, reuse its rollouts, and correct
with weights
\begin{equation}
\omega_j \propto \frac{e^{R_\theta(\tau_j)}}{\pi_\samp(\tau_j)}, \qquad
\widehat{\nabla_\theta L} = \E_{\mathcal{D}_E}[\nabla_\theta R_\theta] -
\sum_j \frac{\omega_j}{\sum_{j'}\omega_{j'}} \nabla_\theta R_\theta(\tau_j),
\label{eq:gclgrad}
\end{equation}
interleaving improvement of $\pi_\samp$ towards the soft optimum. Note the
dynamics terms cancel in the ratio---one never needs $P$---but the weights
degenerate when $\pi_\samp$ drifts from $p_\theta$, which is the practical
failure mode.

\subsection{Support, exploration, and model error}

\begin{itemize}
\item \textbf{Support.} A transition never attempted cannot be estimated.
If $\rho_E$ has support the learner never visits (or vice versa), the
comparison driving \eqref{eq:learningrule} is being made on incomparable
sets. This is the deep reason imitation degrades off the expert's
distribution, and the origin of compounding error in behavioural cloning
\cite{bagnell2014}.
\item \textbf{Model error compounds.} If one fits $\widehat P$ and plans
in it, errors accumulate over the effective horizon: a per-step
total-variation error $\varepsilon$ produces occupancy error of order
$\varepsilon/(1-\gamma)$ and hence a \emph{biased} reward gradient. Bias
does not average away over gradient steps the way variance does.
\item \textbf{Interaction may be unacceptable.} Exploration in surgery,
aviation or finance has costs that no sample-complexity bound captures.
\end{itemize}

\begin{graybox}{A practical decision rule}
Demand exact transition probabilities only when the algorithm asks for
exact dynamic-programming expectations. If you can \emph{sample}
transitions, every such expectation can be estimated without ever writing
$P$ down. What you pay is sample complexity, variance, a coverage
requirement, and possibly interaction risk. What you do \emph{not} pay is
any change to the IRL objective itself: Coordinate 1 is untouched by
Exit C.
\end{graybox}

\subsection{The pivot: from estimating an expectation to estimating a ratio}

Here is the idea that generates the rest of these notes.

The gradient \eqref{eq:learningrule} asks us to compare two expectations
of a \emph{known} function $\phi$ under two \emph{unknown} distributions,
$\rho_E$ and $\rho_\pi$. We have samples from both: the demonstrations
give samples of $\rho_E$, and rollouts give samples of $\rho_\pi$. So the
real task is: \emph{given samples from two distributions, quantify and
reduce their discrepancy.}

Machine learning has an extremely well-developed tool for exactly that
task, and it is not integration. It is \textbf{binary classification}. A
classifier trained to distinguish samples of $\rho_E$ from samples of
$\rho_\pi$ recovers, at its optimum, the density ratio between
them---and the density ratio is precisely the discrepancy signal we
wanted. No partition function, no transition kernel, no enumeration.

That substitution is GAIL. We now make it precise.

\clearpage
% =====================================================================
\part{Adversarial Methods}\label{part:adversarial}
% =====================================================================

\section{Occupancy Measures: the Right Change of Variables}
\label{sec:occvar}

Sections~\ref{sec:selectionprinc1} and \ref{sec:maxent} both optimise over
policies, and both objectives turned out to depend on the policy only
through averages of the form $\sum_{s,a}\rho_\pi(s,a)(\cdot)$. That is a
strong hint that $\rho$, not $\pi$, is the natural variable. Changing to it
converts IRL into a convex-duality problem and makes GAIL a theorem rather
than an analogy.

\begin{definition}[The occupancy polytope]
Let
\begin{equation}
\mathcal{D} \triangleq \Big\{\rho\in\R^{\S\times\A}_{\ge0} \;:\;
\sum_a \rho(s,a) = (1-\gamma)d_0(s) + \gamma\sum_{s',a'} P(s\mid
s',a')\rho(s',a') \;\; \forall s\Big\},
\label{eq:occpolytope}
\end{equation}
the set of non-negative measures satisfying the \emph{Bellman flow}
constraints.
\end{definition}

\begin{theorem}[Policies and occupancies are the same
thing]\label{thm:occbijection}
The map $\pi\mapsto\rho_\pi$ defined by \eqref{eq:occupancydef} is a
bijection from stationary policies (identified when they agree on every
state they visit) onto $\mathcal{D}$, with inverse
\begin{equation}
\pi_\rho(a\mid s) = \frac{\rho(s,a)}{\sum_{a'}\rho(s,a')}.
\label{eq:occinverse}
\end{equation}
$\mathcal{D}$ is a compact convex polytope.
\end{theorem}

This is classical \cite{puterman1994,syed2008}. Three consequences do all
the work:
\begin{enumerate}
\item \textbf{Return is linear in $\rho$:} $J_R(\pi) =
\frac1{1-\gamma}\langle R,\rho_\pi\rangle$ by \eqref{eq:occupancylinear}.
Forward RL is therefore a linear program over $\mathcal{D}$.
\item \textbf{Matching occupancy is matching behaviour}: by
Theorem~\ref{thm:occbijection}, $\rho_\pi=\rho_E$ if and only if
$\pi=\pi_E$ on every state the expert visits. Imitation is exactly a distribution-matching problem, with no
reward mentioned.
\item \textbf{Mixtures are legal}: $\mathcal{D}$ is convex, so
$\rho_{\tilde\pi} = \sum_i \lambda_i \rho_{\pi_i}$ is realisable---which is
what justified the mixed policy in Algorithm~\ref{alg:maxmargin}.
\end{enumerate}

Finally, discounted causal entropy is a concave function of $\rho$:
\begin{equation}
\bar{\mathcal{H}}(\rho) \;=\; -\sum_{s,a} \rho(s,a) \log\frac{\rho(s,a)}{\sum_{a'}\rho(s,a')},
\label{eq:occentropy}
\end{equation}
so ``entropy-regularised RL'' is a concave maximisation over a polytope.

\section{GAIL: Imitation as Density-Ratio Estimation}

\subsection{IRL with an explicit reward regulariser}

Write IRL abstractly, following Ho and Ermon's variational treatment
\cite{hoermon2016}. A reward is good if the expert scores well under it
and the best achievable policy does not score much better---i.e.\ if the
reward \emph{explains the gap}. Add a convex regulariser $\psi$ on rewards
to keep the problem bounded:
\begin{equation}
\mathrm{IRL}_\psi(\pi_E) = \argmax_R \; -\psi(R) \;+\; \Big(\min_\pi
-\beta\bar{\mathcal{H}}(\rho_\pi) - \langle R,\rho_\pi\rangle\Big) \;+\;
\langle R,\rho_E\rangle.
\label{eq:irlpsi}
\end{equation}
Without $\psi$ this is exactly the ill-posed problem of
Section~\ref{sec:invertible}; $\psi$ \emph{is} the selection principle,
written down explicitly for once.

\begin{theorem}[Ho \& Ermon 2016]\label{thm:hoermon}
Running RL on the reward returned by \eqref{eq:irlpsi} yields
\begin{equation}
\mathrm{RL}\circ\mathrm{IRL}_\psi(\pi_E) = \argmin_\pi \; -\beta\bar{\mathcal{H}}(\rho_\pi) \;+\;
\psi^*(\rho_\pi - \rho_E),
\label{eq:rlirl}
\end{equation}
where $\psi^*$ is the convex conjugate of $\psi$.
\end{theorem}

\begin{graybox}{Read Theorem~\ref{thm:hoermon} slowly}
It says that the whole IRL-then-RL pipeline is \emph{equivalent} to a
single occupancy-matching problem, in which the choice of reward
regulariser $\psi$ becomes the choice of \emph{divergence} $\psi^*$ used
to compare $\rho_\pi$ with $\rho_E$. The intermediate reward is a
computational device that can be eliminated. Everything below is a choice
of $\psi$.
\end{graybox}

\subsection{Two choices of $\psi$}

\textbf{(i) A constant on a linear class.} Let $\psi$ be $0$ if $R=w^\top
\phi$ with $\|w\|_2\le1$ and $+\infty$ otherwise. Then $\psi^*$ is the
support function of that set, and \eqref{eq:rlirl} becomes
\begin{equation}
\min_\pi \max_{\|w\|_2\le1} w^\top\big(\mu(\pi)-\mu_E\big) \;=\; \min_\pi
\|\mu(\pi)-\mu_E\|_2,
\label{eq:psilinear}
\end{equation}
which is apprenticeship learning (Section~\ref{sec:selectionprinc1})
recovered exactly. Its weakness is now visible as a property of $\psi^*$:
it is a hard constraint, forcing exact feature matching within a small
class, and it treats any $\rho_\pi$ outside that class as equally bad.

\textbf{(ii) A smooth penalty on the expert's own reward values.} Ho and
Ermon choose $\psi_{\GA}(R) = \E_{\rho_E}[g(R(s,a))]$ for a particular
convex $g$, whose conjugate is a Jensen--Shannon divergence:
\begin{equation}
\psi_{\GA}^*(\rho_\pi-\rho_E) = \max_{D:\S\times\A\to(0,1)}
\E_{\rho_E}\big[\log D(s,a)\big] + \E_{\rho_\pi}\big[\log(1-D(s,a))\big].
\label{eq:psiGA}
\end{equation}
This regulariser penalises rewards that are large on expert pairs, is
defined for \emph{all} rewards rather than a fixed span, and---crucially---its
conjugate is computable from samples alone.

\subsection{The GAIL objective}

\begin{graybox}{GAIL (Ho \& Ermon, 2016)}
\begin{equation}
\min_\pi \max_D \;\; \E_{(s,a)\sim\rho_E}\big[\log D(s,a)\big] +
\E_{(s,a)\sim\rho_\pi}\big[\log\big(1-D(s,a)\big)\big] \;-\;
\beta\bar{\mathcal{H}}(\rho_\pi).
\label{eq:gailobjective}
\end{equation}
\textbf{Convention, fixed for the rest of these notes:} $D(s,a)$ is the
estimated probability that the pair came \textbf{from the expert}. (Ho and
Ermon's paper uses the opposite labelling; the objectives are related by
$D\mapsto1-D$. Mixing the two conventions is the single most common source
of sign confusion in this literature.)
\end{graybox}

\begin{proposition}[The discriminator estimates a density
ratio]\label{prop:gaildratio}
For fixed $\pi$, the inner maximum of \eqref{eq:gailobjective} is attained
at
\begin{equation}
D^\star(s,a) = \frac{\rho_E(s,a)}{\rho_E(s,a)+\rho_\pi(s,a)},
\label{eq:gaildstar}
\end{equation}
and the optimal value equals $2\,\JS(\rho_E\|\rho_\pi) - \log4$. Hence
GAIL is
\begin{equation}
\min_\pi \; 2\,\JS(\rho_E\|\rho_\pi) - \beta\bar{\mathcal{H}}(\rho_\pi).
\label{eq:gailjs}
\end{equation}
\end{proposition}

\begin{proof}
Pointwise in $(s,a)$, maximise $\rho_E\log D + \rho_\pi\log(1-D)$ over
$D\in(0,1)$; the stationary point is $D^\star$. Substituting gives the
Jensen--Shannon expression.
\end{proof}

\subsection{Turning the discriminator into a reward}

The policy player minimises $\E_{\rho_\pi}[\log(1-D)]$, i.e.\ it maximises
the return of the surrogate reward
\begin{equation}
\hat r(s,a) = -\log\big(1-D(s,a)\big) \;\in\; (0,\infty).
\label{eq:gailreward}
\end{equation}
Any policy-gradient method (TRPO, PPO) can then be run on $\hat r$.

\begin{redbox}{Surrogate rewards are not interchangeable}
Three surrogates appear in practice and they are \emph{different
functions}:
\begin{center}
\begin{tabular}{@{}lll@{}}
\toprule
\textbf{surrogate} & \textbf{range} & \textbf{implicit bias}\\
\midrule
$-\log(1-D)$ & $(0,+\infty)$ & always positive: encourages
\emph{surviving} longer\\
$\log D$ & $(-\infty,0)$ & always negative: encourages \emph{terminating}
sooner\\
$\log D - \log(1-D)$ & $(-\infty,\infty)$ & the true log-odds; unbiased in
sign\\
\bottomrule
\end{tabular}
\end{center}
On tasks with variable episode length these choices change the learned
behaviour markedly, independently of anything about imitation. A text that
says ``use the log-odds'' and an equation that uses $-\log(1-D)$ are not
describing the same algorithm. AIRL uses the log-odds form, for reasons
that will be structural rather than cosmetic.
\end{redbox}

\begin{algorithm}[H]
\caption{GAIL (Ho \& Ermon, 2016)}
\begin{algorithmic}[1]
\Require expert data $\mathcal{D}_E$, environment sampler (no $P$
needed), initial $\pi_\vartheta$, discriminator $D_\varphi$
\For{$k=1,2,\ldots$}
  \State collect rollouts $\mathcal{B}\sim\pi_\vartheta$ by interacting
  with the environment
  \State \textbf{discriminator step:} ascend $\nabla_\varphi\big[
  \E_{\mathcal{D}_E}\log D_\varphi + \E_{\mathcal{B}}\log(1-D_\varphi)\big]$
  \State \textbf{policy step:} run one TRPO/PPO update on reward $\hat r =
  -\log(1-D_\varphi)$ with entropy bonus $\beta$
\EndFor
\State \Return $\pi_\vartheta$
\end{algorithmic}
\label{alg:gail}
\end{algorithm}

\subsection{What GAIL does not give you}

\begin{redbox}{At convergence, the discriminator contains no reward
information}
If imitation succeeds then $\rho_\pi=\rho_E$, so $D^\star\equiv\frac12$ and
the surrogate reward \eqref{eq:gailreward} is the \emph{constant} $\log2$.
A constant reward makes every policy optimal
(Section~\ref{sec:invertible}). There is therefore literally nothing to
extract at the end of a successful GAIL run.

Any reward-like signal GAIL exhibits exists only \emph{during} training,
as a by-product of the current mismatch between learner and expert. It is
a function of the optimisation trajectory, not of the task.
This---not a vague appeal to ``shaping terms''---is the precise reason
GAIL does not transfer.
\end{redbox}

\begin{navybox}{Ledger: GAIL}
\textbf{Assumed (Coordinate 1).} Behaviour model: entropy-regularised
optimality. Selection principle: the $\psi_\GA$ regulariser, equivalently
``match occupancy in Jensen--Shannon divergence''. Reward class: whatever
the discriminator architecture can express.\\
\textbf{Approximated (Coordinate 2).} Everything, and cheaply. $Z$ is
never formed; $P$ is never written down; both occupancy measures are
represented only by samples. Requires \emph{online} interaction with the
environment, and inherits the instability of minimax optimisation.\\
\textbf{Delivers.} An imitating policy, often excellent, in continuous
control from raw observations. \textbf{Not} a reward.\\
\textbf{Therefore.} Use GAIL when you want behaviour in the same MDP. If
you want an objective you can re-optimise elsewhere, read on.
\end{navybox}

\section{AIRL: Recovering Something That Transfers}
\label{sec:airl}

AIRL \cite{fu2018} keeps GAIL's sampling machinery and repairs its
Coordinate-1 deficit by \emph{constraining the discriminator} so that a
transferable reward is forced to appear inside it.

\subsection{Where the discriminator's form comes from}

It is not an arbitrary architecture. Suppose the ``expert'' distribution
is the maximum-entropy model $p_\theta(\tau)\propto
q(\tau)e^{\sum_t f_\theta(s_t,a_t)}$ of Section~\ref{sec:maxent}, and the
``fake'' samples come from the current policy $\pi$. Per transition, the
Bayes-optimal discriminator between the two is
\begin{equation}
D_\theta(s,a) = \frac{(\text{expert density})}{(\text{expert
density})+(\text{policy density})} = \frac{\exp f_\theta(s,a)}{\exp
f_\theta(s,a) + \pi(a\mid s)},
\label{eq:airldform}
\end{equation}
because the dynamics factors $q$ are common to both and cancel. Two facts
follow at once.

\begin{proposition}[The AIRL reward is the logit minus a
log-policy]\label{prop:airllogit}
With $D$ as in \eqref{eq:airldform}, the log-odds reward is
\begin{equation}
\hat r(s,a) = \log D_\theta(s,a) - \log\big(1-D_\theta(s,a)\big) =
f_\theta(s,a) - \log\pi(a\mid s).
\label{eq:airlreward}
\end{equation}
\end{proposition}

\begin{remark}[The $-\log\pi$ term is not noise]
Maximising $\E_\pi[\sum_t \gamma^t(f-\log\pi)]$ is exactly maximising
$J_f(\pi) + \mathcal{H}_\gamma(\pi)$: the entropy bonus of
\eqref{eq:softobjective} appears automatically. The adversarial game is
therefore performing entropy-regularised RL on $f$ whether or not anyone
adds an entropy term by hand. At the optimum $f_\theta$ recovers the
\emph{advantage} $A_\soft(s,a) = Q_\soft - V_\soft$ of the expert's soft
policy---which is the crux of what follows, because an advantage is a
reward plus shaping, not a reward.
\end{remark}

\subsection{The structural constraint}

AIRL restricts $f$ to a shaped form:
\begin{equation}
f_{\theta,\omega}(s,a,s') = \underbrace{g_\theta(s)}_{\text{task reward}}
\;+\; \underbrace{\gamma h_\omega(s') - h_\omega(s)}_{\text{potential-based
shaping}}.
\label{eq:airlshaped}
\end{equation}

\begin{redbox}{$g_\theta$ must depend on the state only}
This is the load-bearing assumption of AIRL and it is very often
mis-stated as $g_\theta(s,a)$. Here is why it must be $g_\theta(s)$.

For fixed dynamics $P$, the rewards inducing the same soft-optimal policy
are exactly those of the form
\begin{equation}
r'(s,a) = r(s,a) + \gamma\,\E_{s'\sim P(\cdot\mid s,a)}[\Psi(s')] - \Psi(s)
\label{eq:shapingrepeat}
\end{equation}
for some potential $\Psi$. Note the shaping term $\gamma\E_{s'\mid
s,a}[\Psi(s')]-\Psi(s)$ \textbf{depends on $a$}, and \textbf{depends on
$P$}. So a state--action reward is identified only up to a nuisance term
that (i) cannot be distinguished from the reward by any amount of
behavioural data, and (ii) \emph{changes when the dynamics change}. That
is precisely a failure of transfer.

If instead $g$ is a function of $s$ alone, the only member of the family
\eqref{eq:shapingrepeat} that is state-only is $\Psi\equiv\text{const}$
(under the condition below). The action-dependent freedom is squeezed
entirely into $h_\omega$, and $g_\theta$ is pinned down up to an additive
constant.
\end{redbox}

\begin{definition}[Decomposability condition]
Call two states \emph{one-step linked} if some state has positive-probability
transitions to both. The dynamics are \emph{decomposable} if the
transitive closure of this relation connects all of $\S$. Equivalently:
any two functions $a(s),b(s)$ with $a(s)+b(s')$ constant along all
feasible transitions must themselves be constant.
\end{definition}

\begin{theorem}[Disentanglement, Fu et al.\ 2018]\label{thm:disentangle}
Suppose the dynamics are decomposable, the discriminator is trained to
optimality with $f$ constrained as in \eqref{eq:airlshaped}, and the true
expert is soft-optimal for a state-only reward $r^\star(s)$. Then at the
global optimum there are constants $c_1,c_2$ with
\begin{equation}
g_\theta(s) = r^\star(s) + c_1, \qquad h_\omega(s) = V^\star_\soft(s) + c_2.
\label{eq:disentangleresult}
\end{equation}
In particular $g_\theta$ is invariant to changes in the transition
dynamics, so re-optimising $g_\theta$ in a different MDP with the same
reward structure recovers appropriate behaviour there.
\end{theorem}

\begin{remark}[What the theorem does and does not say]
It says: under a structural condition on the dynamics, a state-only
reward assumption, and exact optimisation, the recovered $g$ equals the
ground-truth reward up to a constant. It does \emph{not} say the reward is
uniquely recoverable in general; the state-only restriction is doing the
identification, and it is an assumption about the world, not a fact
derived from data. If the true reward genuinely depends on actions (an
energy cost, say), AIRL will misattribute that dependence into $h_\omega$
and the recovered $g$ will be wrong. Section~\ref{sec:invertible} has not
been repealed; it has been paid off with an extra assumption, and the
notes' whole point is that this is always the trade.
\end{remark}

\begin{algorithm}[H]
\caption{AIRL (Fu, Luo \& Levine, 2018)}
\begin{algorithmic}[1]
\Require expert data $\mathcal{D}_E$, environment sampler, policy
$\pi_\vartheta$, networks $g_\theta$, $h_\omega$
\For{$k=1,2,\ldots$}
  \State collect rollouts $\mathcal{B}\sim\pi_\vartheta$
  \State form $f_{\theta,\omega}(s,a,s') = g_\theta(s) + \gamma
  h_\omega(s') - h_\omega(s)$ and $D = \dfrac{e^f}{e^f+\pi_\vartheta(a\mid
  s)}$
  \State \textbf{discriminator step:} ascend the binary cross-entropy of
  $D$ on $\mathcal{D}_E$ (label 1) versus $\mathcal{B}$ (label 0), w.r.t.\
  $(\theta,\omega)$
  \State \textbf{policy step:} update $\pi_\vartheta$ by TRPO/PPO on
  $\hat r = \log D - \log(1-D) = f_{\theta,\omega} - \log\pi_\vartheta$
\EndFor
\State \Return transferable reward $g_\theta$; policy $\pi_\vartheta$;
soft value estimate $h_\omega$
\end{algorithmic}
\label{alg:airl}
\end{algorithm}

\begin{navybox}{Ledger: AIRL}
\textbf{Assumed (Coordinate 1).} Everything GAIL assumes, \emph{plus}: the
true reward is state-only, and the dynamics are decomposable. In exchange,
the shaping ambiguity---the one ambiguity that blocks transfer---is
quotiented out by construction.\\
\textbf{Approximated (Coordinate 2).} As GAIL: samples throughout, no $Z$,
no $P$. Additionally requires evaluating $\pi(a\mid s)$ explicitly inside
$D$, so the policy must have a tractable density.\\
\textbf{Delivers.} $g_\theta \approx r^\star + c$, re-optimisable under
new dynamics; $h_\omega \approx V^\star_\soft$, which is a useful
by-product in its own right.\\
\textbf{Fails when.} The reward really depends on actions; the dynamics
are not decomposable; or the discriminator is not trained near
optimality---the theorem is about the global optimum of a minimax problem,
which is not what any optimiser reliably finds.
\end{navybox}

\clearpage
% =====================================================================
\part{Synthesis}\label{part:synthesis}
% =====================================================================

\section{Every Method on the Two Coordinates}

\begingroup
\renewcommand{\arraystretch}{1.2}
\begin{center}
\footnotesize
\textit{Table 2: The same four methods, read as answers to ``what did you
assume?'' and ``what did you approximate?''}\\[4pt]
\begin{tabular}{@{}p{0.19\linewidth}p{0.17\linewidth}p{0.17\linewidth}p{0.17\linewidth}p{0.17\linewidth}@{}}
\toprule
 & \textbf{Max-Margin} & \textbf{MaxCausalEnt} & \textbf{GAIL} & \textbf{AIRL}\\
 & Abbeel \& Ng '04 & Ziebart '08/'10 & Ho \& Ermon '16 & Fu et al.\ '18\\
\midrule
\multicolumn{5}{@{}l}{\textit{Coordinate 1: what is assumed}}\\
\midrule
Behaviour model & expert exactly optimal & Boltzmann-rational, $\pi\propto
e^{Q_\soft/\beta}$ & entropy-regularised optimal & entropy-regularised
optimal\\
Selection principle & largest margin & maximum causal entropy & $\psi_\GA$,
i.e.\ min JS on $\rho$ & $\psi_\GA$ + shaping quotient\\
Reward class & linear in fixed $\phi$ & linear in fixed $\phi$ (or any
$R_\theta$) & implicit: discriminator architecture & $g_\theta(s)$
state-only + $h_\omega$ shaping\\
Extra structural assumption & --- & --- & --- & decomposable dynamics\\
\midrule
\multicolumn{5}{@{}l}{\textit{Coordinate 2: what is approximated}}\\
\midrule
Needs $P$ explicitly? & yes (inner RL + $\mu(\pi)$) & yes:
\eqref{eq:causalQ} and \eqref{eq:forwardrecursion} & no --- needs a
\emph{sampler} & no --- needs a \emph{sampler}\\
Inner expectation via & exact DP / linear solve & exact backward + forward
passes & sampled rollouts + classifier & sampled rollouts + classifier\\
Partition function & absent (no likelihood) & exact, by soft DP & never
formed & never formed\\
Per-iteration cost & QP + full RL solve & $O(T|\S|^2|\A|)$ & rollout batch
+ 2 SGD steps & rollout batch + 2 SGD steps\\
Scales to continuous control? & no & no (tabular) & yes & yes\\
\midrule
\multicolumn{5}{@{}l}{\textit{Outcome}}\\
\midrule
Returns & mixed policy + weak $w$ & unique $w$, soft policy, a likelihood
& policy only & policy + transferable $g_\theta$; $h_\omega \approx
V_\soft$\\
Handles noisy experts & poorly & naturally & naturally & naturally\\
Reward transfers to new dynamics? & only within the class & only within
the class & \textbf{no} ($D\to\frac12$) & yes, under
Thm.~\ref{thm:disentangle}\\
Chief failure mode & infeasible constraints; feature design & intractable
outside small tabular MDPs & minimax instability; no reward at the end &
assumption violation; optimisation never reaches the optimum\\
\bottomrule
\end{tabular}
\end{center}
\endgroup

\begin{graybox}{The one-paragraph version}
Max-margin fixed the ambiguity geometrically and paid with a brittle
behaviour model. Maximum entropy replaced the geometry with a likelihood,
which handled noise and gave a convex problem, and paid with an
intractable partition function. Occupancy measures revealed that both
were solving a distribution-matching problem in disguise. GAIL then
estimated that distributional discrepancy by classification instead of by
integration, which made it scale---at the cost of throwing the reward away
entirely. AIRL put a reward back in by constraining the classifier so that
a state-only component and a shaping component are forced apart, buying
transfer with one more assumption. At no point did anyone recover a true
reward from data alone, and at no point did anyone claim to.
\end{graybox}

\section{The Three Questions, Answered on One Page}

Almost every beginner arrives with the same three questions. Having built
the machinery, we can now answer them crisply.

\begin{graybox}{Q1. ``Everything is finite here---what happens in the
real, infinite case?''}
The \emph{statistics} are unchanged; only the \emph{computation} changes.
Concretely: an infinite horizon needs discounted causal entropy and a soft
Bellman fixed point, which still contracts (Lemma~\ref{lem:lse}), so
nothing is lost. A large or continuous $\S$ destroys the tabular passes
and forces function approximation and sampling. A continuous $\A$ makes
the soft-max over actions an intractable integral, handled variationally.
And on a continuum, the maximum-entropy objective is only well-posed as a
\emph{relative} entropy $-\KL(p\|q)$---which is why we insisted on the
reference measure $q$ in Definition~\ref{def:refmeasure} from the very
first line, even though in the finite deterministic case it is a harmless
constant. The reason to learn IRL discretely first is not that discrete
problems matter; it is that in the discrete case every quantity is
definable, so you can see precisely which ones later become
approximations.
\end{graybox}

\begin{graybox}{Q2. ``Do the features have to be hand-designed?''}
No. The maximum-entropy learning rule is $\E_{\text{expert}}[\nabla_\theta
R_\theta] - \E_{\text{model}}[\nabla_\theta R_\theta]$
(Proposition~\ref{prop:generalgrad}) for \emph{any} differentiable reward
class; the classical rule $\mu_E - \E[\Phi]$ is the special case
$R_\theta=w^\top\phi$. What a fixed linear $\phi$ actually bought was a
finite sufficient statistic, a convex dual, a class-wide performance
guarantee, and interpretability---and what it cost was silent
misspecification. Learning $\phi$ trades convexity for expressiveness and
introduces a new $(w,\phi)\sim(A^{-\top}w,A\phi)$ ambiguity. Most
importantly, the choice never disappears: in GAIL and AIRL the
discriminator architecture \emph{is} the reward class, performed less
visibly.
\end{graybox}

\begin{graybox}{Q3. ``What if we don't know $P$ and have to explore?''}
First separate the two jobs. Estimating the \emph{expert's} statistics
needs no model at all---it is an average over the demonstrations.
Estimating the \emph{current candidate policy's} statistics is the only
place $P$ was ever used in exact MaxEnt IRL: the successor expectation in
the soft backward pass \eqref{eq:causalQ}, and the visitation propagation
in the forward pass \eqref{eq:forwardrecursion}. Both are expectations, so
both can be replaced by sampling whenever you can \emph{draw} transitions,
even if you can never write $P$ down. You then pay in variance, sample
complexity, coverage, and possibly in the real-world risk of exploring;
you do not pay any change in the IRL objective itself. Pushed one step
further, the observation that we only ever needed to \emph{compare} two
sampled distributions replaces integration by classification---and that
step is GAIL.
\end{graybox}

\section{A Checklist for Reading an IRL Paper}

\begin{enumerate}
\item What is the behaviour model? Exact optimality, Boltzmann rationality
with which temperature, or something else?
\item What is the reward class, \emph{including} when it is hidden in an
architecture?
\item What is the selection principle, and is it stated or smuggled?
\item Which quantity plays the role of $\E_{\text{model}}[\cdot]$, and how
is it obtained---exact DP, rollouts, importance weights, or a
discriminator?
\item Does the method need $P$, a simulator, online interaction, or only
offline data?
\item What exactly is claimed to be recovered---a policy, a reward
equivalence class, or a specific reward---and under what conditions?
\item If transfer is claimed, which ambiguity was removed to make transfer
possible, and by what assumption?
\end{enumerate}

\section{Common Errors, Collected}

These are mistakes that recur in secondary expositions, including in
earlier drafts of these notes. Each is discussed at the reference given.

\begin{enumerate}
\item \textbf{Deriving $p(\tau)\propto e^{w^\top\Phi(\tau)}$ and later
asserting $p(\tau)\propto q(\tau)e^{w^\top\Phi(\tau)}$.} These solve
different programs. Put the reference measure into the objective from the
start (Definition~\ref{def:refmeasure}).
\item \textbf{Writing the maximum-entropy program as $-\int p\log p$ on a
continuous trajectory space.} Differential entropy is not
reparameterisation-invariant (Section~\ref{sec:exita}).
\item \textbf{Using the naive soft backward pass under stochastic
dynamics.} It is risk-seeking and implicitly rewrites the transition
kernel (Proposition~\ref{prop:naivetrap}); use maximum causal entropy
(Theorem~\ref{thm:maxcausal}).
\item \textbf{Attributing potential-based shaping to Ng \& Russell
(2000).} It is Ng, Harada \& Russell (1999), and the converse---that this
is the \emph{only} dynamics-independent invariance---is what makes it the
right nuisance class.
\item \textbf{Returning the last policy from Abbeel \& Ng's algorithm.}
The guarantee is for a mixture (Section~\ref{sec:selectionprinc1}).
\item \textbf{Mixing GAIL discriminator conventions, or treating
$-\log(1-D)$, $\log D$ and the log-odds as interchangeable.} They induce
different survival biases (Section~\ref{sec:occvar}).
\item \textbf{Claiming GAIL learns a reward.} At convergence $D\equiv
\frac12$ (Section~\ref{sec:occvar}).
\item \textbf{Writing AIRL's task term as $g_\theta(s,a)$.} The
disentanglement theorem requires a state-only $g$ (Section~\ref{sec:airl}).
\item \textbf{Confusing uniqueness of $V^\star$ with identifiability of
$R$, or uniqueness of the estimator $w$ with identifiability of the
estimand.} (Sections~\ref{sec:invertible} and \ref{sec:maxent}.)
\end{enumerate}

\section{Conclusion}

Inverse reinforcement learning is easy to state and impossible to solve as
stated. The forward map from reward to behaviour destroys information, so
behaviour can never determine reward; and the natural repairs all require
an expectation under the policy that a candidate reward induces, which is
a full forward-RL problem nested inside every gradient step.

Reading the field through those two failures makes its history legible.
Max-margin and maximum entropy are two answers to the first failure, and
maximum entropy wins because it produces a likelihood: noise is modelled
rather than excluded, and the estimation problem becomes convex. The exact
tabular algorithm we built in Section~\ref{sec:maxent} is the reference
implementation of that idea, and its two passes show, with no ambiguity,
where a transition model is required. Every method after it is an answer
to the second failure. Occupancy measures reveal what the inner
expectation really is; sampling replaces the forward pass; importance
weighting reuses samples; and finally the realisation that we only ever
needed to \emph{compare} two sampled distributions replaces integration
with classification, which is GAIL. AIRL then returns to the first
failure with the leverage that adversarial training provides, and buys
back exactly one ambiguity---potential shaping---at the price of exactly
one assumption---a state-only reward.

The lesson worth carrying away is not any single algorithm. It is that in
IRL, every capability is purchased with an assumption, and the
assumptions are always visible if you look for them in the right two
places.

\clearpage
% =====================================================================
\begin{thebibliography}{99}

\bibitem{ngrussell2000}
Ng, A.~Y., \& Russell, S.~J. (2000). Algorithms for inverse reinforcement
learning. In \emph{ICML 2000}, pp.~663--670.

\bibitem{ngharada1999}
Ng, A.~Y., Harada, D., \& Russell, S.~J. (1999). Policy invariance under
reward transformations: theory and application to reward shaping. In
\emph{ICML 1999}, pp.~278--287.

\bibitem{abbeelng2004}
Abbeel, P., \& Ng, A.~Y. (2004). Apprenticeship learning via inverse
reinforcement learning. In \emph{ICML 2004}.

\bibitem{ziebart2008}
Ziebart, B.~D., Maas, A.~L., Bagnell, J.~A., \& Dey, A.~K. (2008). Maximum
entropy inverse reinforcement learning. In \emph{AAAI 2008}, pp.~1433--1438.

\bibitem{ziebart2010thesis}
Ziebart, B.~D. (2010). \emph{Modeling Purposeful Adaptive Behavior with
the Principle of Maximum Causal Entropy}. PhD thesis, Carnegie Mellon
University. (The correct treatment of stochastic dynamics; see
Section~\ref{sec:causaltrap}.)

\bibitem{jaynes1957}
Jaynes, E.~T. (1957). Information theory and statistical mechanics.
\emph{Physical Review}, 106(4), 620--630.

\bibitem{puterman1994}
Puterman, M.~L. (1994). \emph{Markov Decision Processes: Discrete
Stochastic Dynamic Programming}. Wiley. (Occupancy-measure formulation and
the policy--occupancy bijection.)

\bibitem{syed2008}
Syed, U., Bowling, M., \& Schapire, R.~E. (2008). Apprenticeship learning
using linear programming. In \emph{ICML 2008}, pp.~1032--1039.

\bibitem{hoermon2016}
Ho, J., \& Ermon, S. (2016). Generative adversarial imitation learning. In
\emph{NeurIPS 2016}.

\bibitem{fu2018}
Fu, J., Luo, K., \& Levine, S. (2018). Learning robust rewards with
adversarial inverse reinforcement learning. In \emph{ICLR 2018}.

\bibitem{finn2016}
Finn, C., Levine, S., \& Abbeel, P. (2016). Guided cost learning: deep
inverse optimal control via policy optimization. In \emph{ICML 2016},
pp.~49--58.

\bibitem{wulfmeier2015}
Wulfmeier, M., Ondr{\'u}{\v s}ka, P., \& Posner, I. (2015). Maximum entropy
deep inverse reinforcement learning. \emph{arXiv:1507.04888}.

\bibitem{haarnoja2018}
Haarnoja, T., Zhou, A., Abbeel, P., \& Levine, S. (2018). Soft
actor--critic: off-policy maximum entropy deep RL with a stochastic actor.
In \emph{ICML 2018}.

\bibitem{levine2018tutorial}
Levine, S. (2018). Reinforcement learning and control as probabilistic
inference: tutorial and review. \emph{arXiv:1805.00909}.

\bibitem{ramachandran2007}
Ramachandran, D., \& Amir, E. (2007). Bayesian inverse reinforcement
learning. In \emph{IJCAI 2007}, pp.~2586--2591.

\bibitem{cao2021}
Cao, H., Cohen, S., \& Szpruch, L. (2021). Identifiability in inverse
reinforcement learning. In \emph{NeurIPS 2021}.

\bibitem{skalse2023}
Skalse, J., Farnik, L., Motwani, S.~R., Jenner, E., Gleave, A., \& Abate,
A. (2023). Invariance in policy optimisation and partial identifiability
in reward learning. In \emph{ICML 2023}.

\bibitem{watkins1992}
Watkins, C.~J.~C.~H., \& Dayan, P. (1992). Q-learning. \emph{Machine
Learning}, 8(3--4), 279--292.

\bibitem{tsitsiklis1994}
Tsitsiklis, J.~N. (1994). Asynchronous stochastic approximation and
Q-learning. \emph{Machine Learning}, 16(3), 185--202.

\bibitem{sutton1999pg}
Sutton, R.~S., McAllester, D., Singh, S., \& Mansour, Y. (1999). Policy
gradient methods for reinforcement learning with function approximation.
In \emph{NeurIPS 1999}.

\bibitem{azar2013}
Azar, M.~G., Munos, R., \& Kappen, H.~J. (2013). Minimax PAC bounds on the
sample complexity of reinforcement learning with a generative model.
\emph{Machine Learning}, 91(3), 325--349.

\bibitem{suttonbarto2018}
Sutton, R.~S., \& Barto, A.~G. (2018). \emph{Reinforcement Learning: An
Introduction}, 2nd ed. MIT Press.

\bibitem{bagnell2014}
Bagnell, J.~A. (2014). An invitation to imitation. Technical report
CMU-RI-TR-15-08, Carnegie Mellon University.

\end{thebibliography}

\end{document}
